%% 第13回編集委員会（物理第1巻）：組 P1-2
%% p1-04.tex  力-4　高さが速さに変わる：仕事とエネルギー（★4〜5・3段）
%% 数値はすべて Python で計算した：scratchpad/phys1/力-4/calc4.py（g＝9.8 m/s²）
%% 出典（本文で使った実機の値・事実）
%%   ・人が8時間の仕事で続けられる仕事率 約75 W：Marks' Standard Handbook for Mechanical Engineers, 11th ed., p.9-4
%%       （Wikipedia "Human power" の引用で確かめた）
%%   ・軽自動車のエンジンの最高出力 47 kW（64馬力）：自動車業界の自主規制の値
%%   ・地球の GM＝3.986×10^14 m³/s²（WGS 84・IERS の値 3.986004418×10^14 を4けたで）、平均半径 6371 km（巻末の表）
%%   ・月の GM＝4.905×10^12 m³/s²、平均半径 1737.4 km：NASA Moon Fact Sheet（脱出速度 2.38 km/s）
%%   ・ネーターの定理（1918年、エミー・ネーター）
%%   ・ジェームズ・ワット（1736〜1819年）、馬力（約 746 W の英馬力、735.5 W の仏馬力）
\makeatletter
% ベクトル：本文は矢印（第13回の追加の決まり）。太字 \symbf は「専門書をのぞいてみよう」だけ
\newcommand\pxev[1]{\vec{#1}}
% ジェットコースターの車（横から、坂にそって置く）  \pxecar{(位置)}{角度}{色}
\newcommand\pxecar[3]{\begin{scope}[shift={#1}, rotate=#2]
  \filldraw[fill=#3!35, draw=#3!75!black, line width=0.6pt, rounded corners=0.6mm] (-3.6,1.1) rectangle (3.6,4.4);
  \fill[bkInk!80] (-2.2,0.7) circle (0.9) (2.2,0.7) circle (0.9);\end{scope}}
% 台（箱）  \pxebox{(左下)}{はば}{高さ}{色}
\newcommand\pxebox[4]{\filldraw[fill=#4!30, draw=#4!70!black, line width=0.6pt] #1 rectangle ++(#2,#3);}
\makeatother

\phchapter{4}{高さが速さに変わる：\\仕事とエネルギー}{ジェットコースターは、最初の山より高く上がれる？}{\bklevelrange{4}{5}}{高さが速さに変わる：仕事とエネルギー}
\bkmichishirube{4-12,5-12,8-4,8-7,8-10}{力-6,力-13,熱-2}

\begin{goal}
\begin{itemize}[label={\color{bkTheme}●}, leftmargin=1.6\zw, itemsep=0pt, topsep=1pt]
\item \bklevel{3}\enspace 高さと速さの関係を、仕事と運動エネルギーで説明できる。
\item \bklevel{4}\enspace エネルギーの保存と摩擦の損失から、速さ・高さ・仕事率を求められる。
\item \bklevel{5}\enspace $\pxev F=-\mathrm{grad}\,U$ から、位置エネルギーの式を導ける。
\end{itemize}
\end{goal}

\phkoushiki{%
\phk{ゴムをのばすほど、働きが大きい}{小学校理科}{ph:力-4:gomu}
\phk{仕事 $W=Fs$、仕事率、仕事の原理、効率}{中学理科}{ph:力-4:shigoto}
\phk{位置・運動エネルギーとその保存}{中学理科}{ph:力-4:hozon}
\phk{$W=Fs\cos\theta$、$P=Fv$}{高校物理基礎}{ph:力-4:cos}
\phk{$K=\tfrac{1}{2}mv^2$、$U=mgh$、$U=\tfrac{1}{2}kx^2$}{高校物理基礎}{ph:力-4:undou}
\phk{保存力以外の仕事と力学的エネルギー}{高校物理基礎}{ph:力-4:sonshitsu}
\phk{$U=-\tfrac{GMm}{r}$、第2宇宙速度、ループ}{高校物理}{ph:力-4:banyuu}
\phkx{単振動のエネルギー $\tfrac{1}{2}kA^2$}{高校物理}{波-1 で導く}
}

\begin{hatena}
\noindent\begin{minipage}[c]{0.52\linewidth}\raggedright
ジェットコースターは、はじめにチェーンで最初の山の上まで引き上げられ、そこからはモーターなしで走ります。谷へ下りると速くなり、次の山を上ると、おそくなります。

では、2つめの山を、最初の山より高くすることはできるでしょうか。谷でついた速さの勢いで、もっと高くまで上れそうな気もします。
\end{minipage}\hfill
\begin{minipage}[c]{0.45\linewidth}\centering
\begin{tikzpicture}[x=0.82mm, y=0.82mm]
  \draw[bkInk, line width=1.1pt] plot[smooth, tension=0.6] coordinates {(0,26) (6,28) (14,12) (22,3) (30,6) (38,24) (44,32) (50,30)};
  \draw[bkGray, line width=0.4pt, dashed] (0,28) -- (56,28);
  \pxecar{(6,28.2)}{0}{bkRed}
  \node[font=\sffamily\footnotesize, anchor=south, text=bkGray] at (23,28.6) {最初の山の高さ};
  \node[font=\sffamily\bfseries\large, text=bkRed] at (51,37) {?};
  \draw[bkGray, line width=0.6pt] (-2,0) -- (56,0);
\end{tikzpicture}
\end{minipage}
\end{hatena}

\noindent 力-3では、ぶつかったり噴き出したりする動きを、運動量という「数える量」で調べました。この章では、もう1つの数える量、\textbf{エネルギー}を作ります。運動量は力を\textbf{時間}でたし合わせた量でしたが、エネルギーは力を\textbf{距離}でたし合わせた量（仕事）から生まれます。高さと速さが入れかわるしくみがわかると、ジェットコースター・ポンプ・発電所・ロケットまで、同じ考えで計算できます。

\begin{yattemiyou}[ふりこのおもりは、どこまで上がる？]
\textbf{用意するもの}\par
糸（\qty{1}{m}くらい）、おもり（5円玉を数枚）、テープ、えんぴつ、ものさし\par\vspace{2pt}
\noindent\begin{minipage}[c]{0.6\linewidth}
\begin{enumerate}[label={\bkcircled[bkGreen]{\arabic*}}, leftmargin=2.4\zw, itemsep=1pt, topsep=0pt]
\item 糸の先におもりをつけ、かべの前につるして、ふりこを作る。
\item おもりを横へ引き上げ、その高さに印をつけて、そっと放す。反対側で、どの高さまで上がるかを見る。
\item 糸の途中にえんぴつを横に当てて（くぎの代わり）、糸が引っかかってから曲がるようにして、同じ高さから放そう。
\end{enumerate}
\end{minipage}\hfill
\begin{minipage}[c]{0.37\linewidth}\centering
\begin{tikzpicture}[x=0.85mm, y=0.85mm]
  \fill[bkGray!35] (2,40) rectangle (58,42);
  \draw[bkInk, line width=0.6pt] (30,40) -- (30,4);
  \draw[bkGray, line width=0.6pt, dashed] (30,40) -- ($(30,40)+(-130:36)$);
  \fill[bkBlue] ($(30,40)+(-130:36)$) circle (1.6);
  \fill[bkBlue] (30,4) circle (1.6);
  \fill[bkInk] (30,22) circle (0.9);
  \draw[bkGray, line width=0.6pt, dashed] (30,22) -- ($(30,22)+(-32.2:18)$);
  \fill[bkBlue!50] ($(30,22)+(-32.2:18)$) circle (1.6);
  \draw[bkRed, line width=0.6pt, dashed] (4,12.42) -- (56,12.42);
  \node[font=\sffamily\footnotesize, anchor=north, text=bkRed!75!black] at (46,9) {同じ高さ？};
  \node[font=\sffamily\footnotesize, anchor=west] at (31.5,24) {えんぴつ};
\end{tikzpicture}
\end{minipage}
\end{yattemiyou}

\noindent やってみようの答え：おもりは、反対側でも、放したときとほぼ同じ高さまで上がります。えんぴつで糸を途中で曲げても、上がる高さはほぼ同じでした。少しずつ低くなっていくのは、空気の抵抗や、糸の付け根でのこすれがあるからです。どんな道を通っても「もとの高さまでもどろうとする」性質が、この章の主役です（ガリレオが、くぎを打ったふりこで確かめた実験として知られています）。

\begin{omoidasou}
\bkrecall{4-12}{仕事 $W$＝力×力の向きに動いた距離（J）。ななめなら $W=Fs\cos\theta=\pxev F\cdot\pxev s$。重力のする仕事は、下がった高さだけで決まり $mgh$。}
\bkrecall{5-12}{力が変わるときの仕事は $\displaystyle\int F\,\dd x$。ばねをのばす仕事は $\dfrac{1}{2}kx^2$。}
\bkrecall{力-1}{一定の加速度 $a$ で距離 $x$ 進むと、$v^2-v_0^2=2ax$。}
\bkrecall{力-2}{摩擦のない斜面を下る加速度 $g\sin\theta$。動摩擦力 $\mu'N$。}
\bkrecall{8-4・8-7・8-10}{$\mathrm{grad}\,U=\Bigl(\dfrac{\partial U}{\partial x},\dfrac{\partial U}{\partial y},\dfrac{\partial U}{\partial z}\Bigr)$。線積分 $\displaystyle\int_C\pxev F\cdot\dd\pxev r$ は、道にそって「力×小さな移動」をたす。$\mathrm{rot}\,\pxev F=\pxev 0$ なら、線積分は道によらない。}
\end{omoidasou}

%% ======================== 物理の道具 ========================
\bksection{doguu}{物理の道具}

\phstage{me}{高さと速さの関係を見る}{3}

\bksubsection{1}{コースの高さと速さ}
\label{ph:力-4:sec1}
\noindent\begin{minipage}[c]{0.5\linewidth}\centering
\begin{tikzpicture}[x=0.85mm, y=0.85mm]
  \draw[bkInk, line width=1.1pt] plot[smooth, tension=0.5] coordinates {(0,40) (6,40.5) (14,30) (22,20) (30,10) (38,1) (46,0) (56,0)};
  \foreach \yy in {0,10,20,30,40} {\draw[bkGray, line width=0.3pt, dashed] (-3,\yy) -- (60,\yy); \node[font=\sffamily\footnotesize, anchor=east] at (-3.5,\yy) {\yy\,m};}
  \pxecar{(6,40.6)}{0}{bkRed}
  \foreach \px/\py/\vv in {14/30/14.0, 22/20/19.8, 30/10/24.2, 50/0/28.0} {\fill[bkBlue] (\px,\py) circle (0.9); \node[font=\sffamily\footnotesize, text=bkBlue!80!black, anchor=south west, fill=white, inner sep=0.8pt] at (\px+1,\py+1.2) {\vv};}
  \node[font=\sffamily\footnotesize, anchor=south, text=bkRed!75!black] at (8,46) {速さ0};
\end{tikzpicture}
\end{minipage}\hfill
\begin{minipage}[c]{0.47\linewidth}\raggedright
高さ\qty{40}{m}の山の上から、速さ0で下り始めるコースターの速さ（m/s。青い数）を、摩擦と空気の抵抗がないとして計算しました（どう計算したかは③で）。\par\vspace{2pt}
\qty{10}{m}下りるたびに速くなりますが、速さのふえ方は、だんだん小さくなっています。
\end{minipage}

\par\vspace{4pt}
速さのかわりに、速さの2乗 $v^2$ を表にしてみましょう。
\begin{bkcenter}
\begin{bktable}{colspec={Q[l,m,wd=40mm] *{5}{Q[c,m,wd=14mm]}}, rowsep=1.5pt, row{1}={bg=bkTheme!10}}
高さ（m） & 40 & 30 & 20 & 10 & 0 \\
下りた高さ（m） & 0 & 10 & 20 & 30 & 40 \\
速さ $v$（m/s） & 0 & 14.0 & 19.8 & 24.2 & 28.0 \\
$v^2$（m$^2$/s$^2$） & 0 & 196 & 392 & 588 & 784 \\
\end{bktable}
\bktabcaption{速さの2乗は、下りた高さに比例する（\qty{10}{m}ごとに196ずつふえる）。}
\end{bkcenter}
$v^2$ は、\qty{10}{m}下りるごとにちょうど196ずつふえています。1 m あたりでは $196\div10=19.6$ で、これは $2\times9.8$ です。だから、
\[v^2=2g\times(\text{下りた高さ})\]
です。これは、力-1の自由落下の式 $v^2=2gh$ と同じです。坂がゆるくても急でも、まっすぐ落ちても、\textbf{下りた高さが同じなら、速さは同じ}になるようです。なぜでしょう。

\bksubsection{2}{仕事：力×力の向きに動いた距離}
\label{ph:力-4:sec2}
4-12で、仕事を「力×力の向きに動いた距離」と決めました。重力が物にする仕事は、坂の角や長さによらず、下がった高さ $h$ だけで決まり、$mgh$ でした。

\noindent\begin{minipage}[c]{0.5\linewidth}\centering
\begin{tikzpicture}[x=1mm, y=1mm]
  \draw[bkGray, line width=0.6pt] (-2,0) -- (62,0);
  \draw[bkGray, line width=0.4pt, dashed] (-2,24) -- (62,24);
  \draw[bkBlue, line width=1.3pt] (4,24) -- (4,0);
  \draw[bkOrange, line width=1.3pt] (14,24) -- (28,0);
  \draw[bkGreen!70!black, line width=1.3pt] (34,24) -- (62,0);
  \fill[bkRed] (4,24) circle (1.2) (14,24) circle (1.2) (34,24) circle (1.2);
  \draw[<->, bkInk, line width=0.6pt] (-1,0.5) -- (-1,23.5);
  \node[font=\sffamily\footnotesize, anchor=east] at (-1.5,12) {$h$};
  \node[font=\sffamily\footnotesize, anchor=north] at (31,-1.5) {どの道でも、重力の仕事は $mgh$};
\end{tikzpicture}
\end{minipage}\hfill
\begin{minipage}[c]{0.47\linewidth}\raggedright
急な坂は短く、ゆるい坂は長い。坂にそった重力の成分 $mg\sin\theta$ と、坂の長さ $\dfrac{h}{\sin\theta}$ をかけると、$\sin\theta$ が消えて、いつも $mgh$ です。
\end{minipage}

\par\vspace{4pt}
①の表では、速さの2乗 $v^2$ が、下がった高さ（＝重力のした仕事）に比例していました。どうやら、\textbf{重力のした仕事が、速さの2乗に変わっている}ようです。このことを、③で数で確かめます。

\Needspace{16\baselineskip}\phstage{te}{表計算で、坂を区切って追う}{3}

\bksubsection{3}{曲がった坂を、短いまっすぐな坂に分ける}
\label{ph:力-4:sec3}
曲がったコースは、短いまっすぐな坂をつないだもの、と考えます。1つのまっすぐな坂（角 $\theta$、長さ $\Delta s$）の上では、加速度は一定です。

\noindent\begin{minipage}[c]{0.46\linewidth}\centering
\begin{tikzpicture}[x=1mm, y=1mm]
  \draw[bkInk, line width=1.1pt] (0,24) -- (44,4);
  \draw[bkGray, line width=0.5pt, dashed] (0,4) -- (48,4) (0,24) -- (0,4);
  \draw[<->, bkBlue, line width=0.6pt] (0,1) -- (44,1);
  \node[font=\sffamily\footnotesize, anchor=north, text=bkBlue!80!black] at (22,0.5) {$\Delta x$（水平の距離）};
  \draw[<->, bkOrange, line width=0.6pt] (-2,4.5) -- (-2,23.5);
  \node[font=\sffamily\footnotesize, anchor=east, text=bkOrange!75!black] at (-2.5,14) {$-\Delta h$};
  \node[font=\sffamily\footnotesize, anchor=south west] at (22,15) {$\Delta s$};
  \draw[bkInk, line width=0.5pt] (36,4) arc (180:155.6:8);
  \node[font=\sffamily\footnotesize] at (33,6.8) {$\theta$};
  \pxecar{(16,16.7)}{-24.4}{bkRed}
\end{tikzpicture}
\end{minipage}\hfill
\begin{minipage}[c]{0.51\linewidth}\raggedright
坂を下る向きの加速度は、力-2から
\[a=g\sin\theta-\mu'g\cos\theta\]
です（$\mu'$ は動摩擦係数。転がる抵抗などをまとめて表す）。\par\vspace{2pt}
下がった高さは $-\Delta h=\Delta s\sin\theta$、水平に進んだ距離は $\Delta x=\Delta s\cos\theta$ です。
\end{minipage}

\par\vspace{4pt}
力-1の $v^2-v_0^2=2ax$ で、進んだ距離 $x$ を坂にそった距離 $\Delta s$ にして、$a$ を入れると、
\[v^2-v_0^2=2g\,\Delta s\sin\theta-2\mu'g\,\Delta s\cos\theta=2g(-\Delta h)-2\mu'g\,\Delta x\]
$\theta$ が消えて、\textbf{高さの変化 $\Delta h$ と水平の距離 $\Delta x$ だけ}で、次の点の $v^2$ が決まりました。上り坂でも、加速度が $-g\sin\theta-\mu'g\cos\theta$（摩擦はいつも動きと反対向き）になり、$\Delta h$ が正になるだけで、同じ式が使えます。これをくり返すだけで、どんな形のコースでも計算できます。

\phboxhead{表計算で追う}
\qty{10}{m}ごとの高さを決めたコースで、摩擦なし（$\mu'=0$）と、摩擦あり（$\mu'=0.03$）の $v^2$ を計算します。
\begin{bkcenter}
\bksheet{6}{17}{%
1 & $x$（m） & 高さ（m） & $v^2$（摩擦なし） & $v$（m/s） & $v^2$（摩擦あり） & $v$（m/s） \\
2 & 0 & 40 & 0.0 & 0.0 & 0.0 & 0.0 \\
3 & 10 & 34 & 117.6 & 10.8 & 111.7 & 10.6 \\
4 & 20 & 22 & 352.8 & 18.8 & 341.0 & 18.5 \\
5 & 30 & 10 & 588.0 & 24.2 & 570.4 & 23.9 \\
6 & 40 & 3 & 725.2 & 26.9 & 701.7 & 26.5 \\
7 & 50 & 2 & 744.8 & 27.3 & 715.4 & 26.7 \\
8 & 60 & 8 & 627.2 & 25.0 & 591.9 & 24.3 \\
9 & 70 & 18 & 431.2 & 20.8 & 390.0 & 19.7 \\
10 & 80 & 28 & 235.2 & 15.3 & 188.2 & 13.7 \\
11 & 90 & 35 & 98.0 & 9.9 & 45.1 & 6.7 \\
12 & 100 & 38 & 39.2 & 6.3 & $-19.6$ & ― \\}
\bktabcaption{C3＝C2－2*9.8*(B3－B2)、E3＝E2－2*9.8*(B3－B2)－2*0.03*9.8*(A3－A2)、D3＝SQRT(C3)。}
\end{bkcenter}

\begin{bkfigure}
\begin{tikzpicture}
\begin{axis}[bkplot, name=tr, width=120mm, height=40mm, xmin=0, xmax=108, ymin=0, ymax=52, xtick={20,40,60,80,100}, ytick={20,40}, xlabel={$x$（m）}, ylabel={高さ（m）}]
  \draw[bkInk, line width=1.3pt] plot[smooth, tension=0.4] coordinates {(axis cs:0,40) (axis cs:10,34) (axis cs:20,22) (axis cs:30,10) (axis cs:40,3) (axis cs:50,2) (axis cs:60,8) (axis cs:70,18) (axis cs:80,28) (axis cs:90,35) (axis cs:100,38)};
  \addplot[bkRed, line width=0.8pt, dashed, domain=0:104] {40};
  \fill[bkOrange] (axis cs:97,37.1) circle (2pt);
  \node[font=\sffamily\footnotesize, text=bkOrange!75!black, anchor=south east] at (axis cs:99,41.5) {摩擦ありはここで止まる};
  \node[font=\sffamily\footnotesize, text=bkRed!75!black, anchor=south west] at (axis cs:3,41.5) {最初の山の高さ 40 m};
\end{axis}
\end{tikzpicture}
\bkfigcaption{表計算のコース。2つめの山（\qty{38}{m}）は最初の山より低いのに、摩擦があると上りきれない。}\label{ph:力-4:fig-course}
\end{bkfigure}

\phboxhead{表から読みとれること}
\begin{enumerate}[label={\bkcircled{\arabic*}}, leftmargin=2.4\zw, itemsep=1pt, topsep=2pt]
\item 摩擦なしでは、$v^2$ はいつも $2g\times(40-\text{高さ})$ です。たとえば高さ\qty{38}{m}で $2\times9.8\times2=39.2$。坂の形によりません。
\item 摩擦ありでは、水平に\qty{10}{m}進むごとに、$v^2$ が摩擦なしより5.88ずつ余分に減ります（$2\times0.03\times9.8\times10$）。\qty{90}{m}進むと、52.9も減っています（$5.88\times9$）。
\item だから、摩擦なしでは、2つめの山（\qty{38}{m}）を速さ\qty{6.3}{m/s}で越えますが、摩擦ありでは $v^2$ が負になり、上りきれません（$x$ が約\qty{97}{m}、高さ約\qty{37}{m}で止まり、もどってきます）。
\end{enumerate}

\phboxhead{両辺に $\dfrac{1}{2}m$ をかけると：運動エネルギーと仕事}
上の式 $v^2-v_0^2=2g(-\Delta h)-2\mu'g\,\Delta x$ の両辺に、$\dfrac{1}{2}m$ をかけてみます。
\[\frac{1}{2}mv^2-\frac{1}{2}mv_0^2=\underbrace{mg\times(-\Delta h)}_{\text{重力がした仕事}}+\underbrace{(-\mu'mg\cos\theta)\times\Delta s}_{\text{摩擦力がした仕事}}\]
右辺の1つめは、重力が下がった高さの分だけした仕事（4-12）、2つめは、大きさ $\mu'mg\cos\theta$ の摩擦力が、動く向きと反対向きにした（負の）仕事です。左辺の $\dfrac{1}{2}mv^2$ を\textbf{運動エネルギー}（$K$）とよぶと、
\[\text{運動エネルギーの変化}=\text{物がされた仕事の和}\]
と言えます。重力は、仕事をして運動エネルギーをふやし、摩擦は、負の仕事をして運動エネルギーを減らします。\bkidx{うんどうえねるぎー}{運動エネルギー}

\phboxhead{はてな？の答え（\qty{40}{m}の山から始めるとき）}
摩擦のないコースなら、コースターはちょうど最初の山と同じ高さまでしか上がれません（$v^2=0$ になる所）。本物のコースには、車輪の転がる抵抗と空気の抵抗があるので、最初の山と同じ高さにも届きません。だから、\textbf{モーターやチェーンで引き上げないかぎり、2つめの山は最初の山より低く作ります}。⑦で、式でも確かめます。

\begin{phmodel}[坂を区切った計算]
\phignore{コースの曲がり（遠心力の分の垂直抗力の変化）。空気の抵抗の、速さによる変化。車輪の回転のエネルギー。}
\phassume{摩擦力は $\mu'mg\cos\theta$ で、$\mu'$ は一定。区切った1つの坂の上で加速度は一定。}
\phlimit{区切りが粗すぎるとき。ループや急なカーブで、垂直抗力が $mg\cos\theta$ から大きくずれるとき（⑨）。速くて空気の抵抗が大きいとき。}
\end{phmodel}
\phkokomade

\Needspace{16\baselineskip}\phstage{shiki}{線積分とポテンシャルで導く}{5}

\bksubsection{4}{仕事を線積分で書く}
\label{ph:力-4:sec4}
③では、坂を短いまっすぐな坂に分けて、それぞれの「力×力の向きに動いた距離」をたしました。これを、どんな道・どんな力にも使える形にします。

道を小さな移動 $\Delta\pxev r_i$ に分けます。それぞれの区間で、力 $\pxev F_i$ はほぼ一定なので、その区間の仕事は内積 $\pxev F_i\cdot\Delta\pxev r_i$（4-12）です。それをたして、区切りを細かくした先が\textbf{線積分}（8-7）です。
\pheq{W=\int_C\pxev F\cdot\dd\pxev r=\lim\sum_i\pxev F_i\cdot\Delta\pxev r_i}{N·m ＝ J}
$C$ は、物が通った道を表します。これが、仕事のいちばん一般の形です。

\begin{phflow}[仕事]
\phflowrow{$W=\displaystyle\int_C\pxev F\cdot\dd\pxev r$}{$\pxev F$ が一定\\道がまっすぐ}{$W=\pxev F\cdot\Delta\pxev r=Fs\cos\theta$}
\phflowrow{$W=Fs\cos\theta$}{力と動く向きが\\同じ（$\theta=0$）}{$W=Fs$}
\end{phflow}
\label{ph:力-4:cos}力が一定で、道がまっすぐなら、どの区間も同じ $\pxev F\cdot\Delta\pxev r_i$ なので、和は $\pxev F\cdot\Delta\pxev r$、大きさと角で書くと $Fs\cos\theta$ です。\label{ph:力-4:shigoto}さらに力と動く向きが同じなら、$\cos0^\circ=1$ で $W=Fs$。中学校の「仕事＝力×力の向きに動いた距離」です。

\begin{phyomi}[$W=\displaystyle\int_C\pxev F\cdot\dd\pxev r$]
「通った道にそって、少し進むごとに、力の『進む向きの成分』×進んだ距離をたしていく」。力が動く向きと直角なら、その区間の仕事は0（垂直抗力・向心力は仕事をしない。4-12）。
\end{phyomi}

\begin{phbaisuru}[いろいろな角]{力と動く向きの角 $\theta$}
$F=\qty{10}{N}$、$s=\qty{2}{m}$ なら、$\theta=0^\circ$ で\qty{20}{J}、$60^\circ$ で\qty{10}{J}、$90^\circ$ で0、$180^\circ$ で $\qty{-20}{J}$（4-12の表）。動きをじゃまする向きの力は、負の仕事をします。
\end{phbaisuru}

%% ---------------------------------------------------------------
\bksubsection{5}{運動エネルギーと仕事}
\label{ph:力-4:sec5}
運動方程式 $m\dfrac{\dd\pxev v}{\dd t}=\pxev F$ の両辺と、速度 $\pxev v$ との内積をとります。
\[m\,\pxev v\cdot\frac{\dd\pxev v}{\dd t}=\pxev F\cdot\pxev v\]
左辺は、積の微分（5-6）の
\[\frac{\dd}{\dd t}\bigl(\pxev v\cdot\pxev v\bigr)=2\,\pxev v\cdot\frac{\dd\pxev v}{\dd t}\]
を使うと、$\dfrac{\dd}{\dd t}\Bigl(\dfrac{1}{2}mv^2\Bigr)$ と書けます（$\pxev v\cdot\pxev v=v^2$）。右辺は、$\pxev v\,\dd t=\dd\pxev r$（少しの時間に動いた分）なので、時間で積分すると線積分になります。時刻 $t_1$ から $t_2$ まで積分して、
\pheq{\frac{1}{2}mv_2^{\,2}-\frac{1}{2}mv_1^{\,2}=\int_C\pxev F\cdot\dd\pxev r}{kg·m$^2$/s$^2$ ＝ N·m ＝ J}
\label{ph:力-4:undou}左辺の $K=\dfrac{1}{2}mv^2$ が\textbf{運動エネルギー}です。「運動エネルギーの変化は、物がされた仕事（すべての力の仕事の和）に等しい」。③で、まっすぐな坂ごとに確かめたことが、どんな道・どんな力でも成り立つとわかりました。

\begin{phyomi}[$K=\dfrac{1}{2}mv^2$]
「動いている物がもつ、仕事をする能力」。$\dfrac{1}{2}$ は、速さが0から $v$ まで一様にふえるあいだの平均の速さ $\dfrac{v}{2}$ から来る（5-12のばねの $\dfrac{1}{2}$ と同じ三角形の面積）。単位は J。
\end{phyomi}

\begin{phmijika}[止まるまでの距離]
ブレーキで、一定の力 $F$ で止めると、止まるまでの距離 $d$ は、$Fd=\dfrac{1}{2}mv^2$ から $d=\dfrac{mv^2}{2F}$ です。力-1で求めた $\dfrac{v^2}{2a}$ と同じ式ですが、エネルギーで考えると、「運動エネルギーを、ブレーキの仕事で0にする」と一言で言えます。
\end{phmijika}

\begin{phbaisuru}{速さ $v$}
運動エネルギーは $v^2$ に比例するので、4倍になります。同じブレーキで止まるまでの距離も4倍です。時速\qty{40}{km}から時速\qty{80}{km}にすると、ブレーキをかけてから止まるまでの距離は4倍になります。
\par\smallskip\noindent\begin{tikzpicture}
\begin{axis}[bkplot, width=96mm, height=42mm, xmin=0, xmax=33, ymin=0, ymax=520, xtick={10,20,30}, ytick={50,200,450}, xlabel={$v$（m/s）}, ylabel={$K$（J）}]
  \addplot[bkcurve, domain=0:31] {0.5*x^2};
  \fill[bkRed] (axis cs:10,50) circle (1.8pt) (axis cs:20,200) circle (1.8pt) (axis cs:30,450) circle (1.8pt);
  \node[font=\sffamily\footnotesize, anchor=north west, text=bkRed!75!black] at (axis cs:20.5,195) {2倍の速さで4倍};
  \node[font=\sffamily\footnotesize, anchor=north west, text=bkGray] at (axis cs:1,500) {質量 \qty{1}{kg} のとき};
\end{axis}
\end{tikzpicture}
\end{phbaisuru}

\begin{phflow}[運動エネルギー]
\phflowrow{$m\dfrac{\dd\pxev v}{\dd t}=\pxev F$}{$\pxev v$ と内積をとり\\時間で積分}{$K_2-K_1=\displaystyle\int_C\pxev F\cdot\dd\pxev r$}
\phflowrow{$K_2-K_1=\displaystyle\int_C\pxev F\cdot\dd\pxev r$}{一直線上・\\$F$ が一定}{$\dfrac{1}{2}mv'^2-\dfrac{1}{2}mv^2=Fs$}
\end{phflow}

%% ---------------------------------------------------------------
\bksubsection{6}{保存力とポテンシャル}
\label{ph:力-4:sec6}
\phboxhead{道によらない仕事}
②で、重力のする仕事は、道によらず高さの差だけで決まりました。このように、\textbf{はじめの点とおわりの点だけで仕事が決まる力}を\textbf{保存力}といいます。\bkidx{ほぞんりょく}{保存力}保存力なら、点 $\pxev r$ ごとに1つの数 $U(\pxev r)$ を決めて、
\pheq{\int_{\mathrm A}^{\mathrm B}\pxev F\cdot\dd\pxev r=U(\mathrm A)-U(\mathrm B)}{J ＝ J}
と書けます。この $U$ を\textbf{位置エネルギー}（ポテンシャルエネルギー）といいます。\bkidx{いちえねるぎー}{位置エネルギー}「力がする仕事＝位置エネルギーの減った分」です。小さな移動 $\dd\pxev r$ で書くと $\pxev F\cdot\dd\pxev r=-\dd U$ で、これは8-4の勾配を使って
\pheq{\pxev F=-\mathrm{grad}\,U=-\Bigl(\frac{\partial U}{\partial x},\ \frac{\partial U}{\partial y},\ \frac{\partial U}{\partial z}\Bigr)}{N ＝ J/m}
と同じことです。\textbf{力は、位置エネルギーが下がる向きに、下がり方が急なほど強くはたらきます}。地図の等高線と坂の関係（8-4：ボールは等高線と直角に、急な向きへ転がる）と同じです。

\phboxhead{重力：$U=mgh$}
地面の近くで、上向きに $z$ 軸をとると、重力は $\pxev F=(0,\ 0,\ -mg)$ です。$-\dfrac{\partial U}{\partial z}=-mg$、$\dfrac{\partial U}{\partial x}=\dfrac{\partial U}{\partial y}=0$ になる関数は、$U=mgz$（＋定数）です。高さ $h$ の所なら
\pheq{U=mgh}{kg × m/s$^2$ × m ＝ J}
\label{ph:力-4:mgh}です。これが\textbf{重力による位置エネルギー}です。どこを高さ0にするかは自由に決めてよく（定数の分）、仕事に効くのは差だけです。

\begin{phbaisuru}{高さ $h$}
$U=mgh$ は $h$ に比例して2倍になります。でも、そこから落ちたときの速さは、$\dfrac{1}{2}mv^2=mgh$ より $v=\sqrt{2gh}$ なので、$\sqrt2\approx1.41$ 倍にしかなりません。速さを2倍にするには、高さを4倍にする必要があります。
\par\smallskip\noindent\begin{tikzpicture}
\begin{axis}[bkplot, width=96mm, height=40mm, xmin=0, xmax=85, ymin=0, ymax=45, xtick={20,40,60,80}, ytick={14,19.8,28,39.6}, yticklabels={14,19.8,28,39.6}, xlabel={$h$（m）}, ylabel={$v$（m/s）}]
  \addplot[bkcurve, domain=0:82, samples=100] {sqrt(2*9.8*x)};
  \fill[bkRed] (axis cs:10,14) circle (1.8pt) (axis cs:20,19.8) circle (1.8pt) (axis cs:40,28) circle (1.8pt) (axis cs:80,39.6) circle (1.8pt);
  \node[font=\sffamily\footnotesize, anchor=north west, text=bkRed!75!black] at (axis cs:41,26.5) {高さ4倍で速さ2倍};
\end{axis}
\end{tikzpicture}
\end{phbaisuru}

\phboxhead{ばね：$U=\dfrac{1}{2}kx^2$}
自然の長さからののびを $x$ とすると、ばねの力は $F=-kx$（のびと反対向き。力-2）です。$-\dfrac{\dd U}{\dd x}=-kx$ になる関数は
\pheq{U=\frac{1}{2}kx^2}{N/m × m$^2$ ＝ J}
\label{ph:力-4:bane}で、5-12で「ばねをのばす仕事」として三角形の面積から求めた値と同じです。これが\textbf{弾性力による位置エネルギー}（弾性エネルギー）です。\label{ph:力-4:gomu}小学校の理科で調べた「ゴムを長くのばすほど、物を動かす働きが大きい」は、このことです。ゴムを2倍のばすと、たくわえられるエネルギーは（ゴムの力がのびに比例する範囲なら）4倍になります。

\begin{phbaisuru}{のび $x$}
$U=\dfrac{1}{2}kx^2$ は $x^2$ に比例するので4倍です。おもちゃの鉄砲のばねを2倍縮めると、玉の運動エネルギーは4倍、速さは2倍になります（例題4）。
\end{phbaisuru}

\phboxhead{万有引力：$U=-\dfrac{GMm}{r}$}
地球（質量 $M$）の中心から距離 $r$ の所で、質量 $m$ の物が受ける万有引力は、中心向きに大きさ $\dfrac{GMm}{r^2}$ です（力-2）。中心から外向きの単位ベクトルを $\pxev e_r$ とすると $\pxev F=-\dfrac{GMm}{r^2}\pxev e_r$ です。$r$ の向きに微分して $-\dfrac{\dd U}{\dd r}=-\dfrac{GMm}{r^2}$ になる関数は、$\dfrac{\dd}{\dd r}\Bigl(-\dfrac{1}{r}\Bigr)=\dfrac{1}{r^2}$ から
\pheq{U=-\frac{GMm}{r}}{N·m$^2$/kg$^2$ × kg$^2$/m ＝ J}
\label{ph:力-4:banyuu}です（無限に遠い所を0とした）。負の数なのは、無限に遠い所より「下」にいる、という意味です。

\begin{bkfigure}
\begin{tikzpicture}
\begin{axis}[bkplot, width=112mm, height=50mm, xmin=0, xmax=5.3, ymin=-1.15, ymax=0.25, xtick={1,2,3,4,5}, xticklabels={$R$,$2R$,$3R$,$4R$,$5R$}, ytick={-1,-0.5}, yticklabels={$-\frac{GMm}{R}$,}, xlabel={$r$}, ylabel={$U$}, axis x line=middle]
  \addplot[bkcurve, domain=1:5.1, samples=80] {-1/x};
  \addplot[bkcurve2, domain=1:1.6] {-1+(x-1)};
  \fill[bkGray!30] (axis cs:0,-1.15) rectangle (axis cs:1,0);
  \node[font=\sffamily\footnotesize, anchor=west, text=bkOrange!75!black] at (axis cs:1.62,-0.42) {地面の近くでは $mgh$（直線）};
  \node[font=\sffamily\footnotesize, anchor=west, text=bkBlue!80!black] at (axis cs:2.6,-0.24) {$U=-\dfrac{GMm}{r}$};
  \node[font=\sffamily\footnotesize] at (axis cs:0.5,-0.6) {地球};
\end{axis}
\end{tikzpicture}
\bkfigcaption{万有引力の位置エネルギー。地球の表面の近くでは、曲線をほぼ直線（傾き $mg$）とみなせる。}\label{ph:力-4:fig-banyuu}
\end{bkfigure}

\phboxhead{地面の近くでは $mgh$ にもどる}
地球の半径を $R$、地面からの高さを $h$ とすると $r=R+h$ です。$h$ が $R$ よりずっと小さい（$h\ll R$）とき、$\dfrac{1}{1+x}\approx1-x$（5-17、$x$ が小さいとき）を使って、
\[U=-\frac{GMm}{R+h}=-\frac{GMm}{R}\cdot\frac{1}{1+\dfrac{h}{R}}\approx-\frac{GMm}{R}+\frac{GMm}{R^2}h\]
です。力-2で $g=\dfrac{GM}{R^2}$ だったので、2つめの項は $mgh$ です。1つめの項は $h$ によらない定数なので、差だけを考えるときには消えます。だから、地面の近くでは $U=mgh$ と同じです。

\begin{phflow}[位置エネルギー]
\phflowrow{$\pxev F=-\mathrm{grad}\,U$}{重力 $(0,\,0,\,-mg)$\\地面の近く}{$U=mgh$}
\phflowrow{$\pxev F=-\mathrm{grad}\,U$}{ばねの力\\$F=-kx$}{$U=\dfrac{1}{2}kx^2$}
\phflowrow{$U=-\dfrac{GMm}{r}$}{$r=R+h$\\$h\ll R$}{$U=mgh$（＋定数）}
\end{phflow}

\phboxhead{地球から飛び出す速さ：第2宇宙速度}
\label{ph:力-4:uchuu}地面から真上に投げた物が、地球にもどってこないで無限に遠くまで行くための、いちばん小さい速さを求めます。空気の抵抗を無視すると、運動エネルギーと万有引力の位置エネルギーの和は変わりません（⑦）。無限に遠い所で速さがちょうど0になるとして、
\[\frac{1}{2}mv^2-\frac{GMm}{R}=0+0\qquad v=\sqrt{\frac{2GM}{R}}=\sqrt{2gR}\]
この速さを\textbf{第2宇宙速度}といいます。地球の値（$GM$ が\qty{3.986e14}{m^3/s^2}、$R$ が\qty{6.371e6}{m}）を入れると、約\qty{11.2}{km/s}です（例題11）。人工衛星の速さ（第1宇宙速度）$\sqrt{gR}$ の $\sqrt2$ 倍です。

\phboxhead{保存力の見分け方：$\mathrm{rot}\,\pxev F=\pxev 0$}
どんな力でも位置エネルギーが作れるわけではありません。8-10で学んだように、$\pxev F=-\mathrm{grad}\,U$ と書ける力では $\mathrm{rot}\,\pxev F=\pxev 0$ になり、反対に（すきまのない広がりの中なら）$\mathrm{rot}\,\pxev F=\pxev 0$ の力は $-\mathrm{grad}\,U$ と書けます。平面の力 $\pxev F=(F_x,\ F_y)$ なら、確かめるのは
\[\frac{\partial F_y}{\partial x}-\frac{\partial F_x}{\partial y}=0\ \text{かどうか}\]
だけです。たとえば $\pxev F=(y,\ x)$ は $1-1=0$ で保存力（$U=-xy$）、$\pxev F=(-y,\ x)$ は $1-(-1)=2$ で保存力ではありません（例題13）。

\begin{phcode}
from math import cos, sin, pi
N = 100                  \# 円を N 本の線分に分ける
W = 0.0
for i in range(N):
    t0, t1 = 2*pi*i/N, 2*pi*(i+1)/N
    x0, y0 = cos(t0), sin(t0)
    x1, y1 = cos(t1), sin(t1)
    xm, ym = (x0+x1)/2, (y0+y1)/2  \# まん中の点
    W = W + (-ym)*(x1-x0) + xm*(y1-y0)
print(round(W, 4))
\end{phcode}
\begin{phcodeout}
6.2791
\end{phcodeout}
W をふやす行は、力 $(-y,\ x)$ と小さな移動 $(x_1-x_0,\ y_1-y_0)$ の内積です。半径1の円を1周する仕事は、分け方を細かくすると $2\pi\approx6.283$ に近づき、0になりません。同じ点にもどってきたのに仕事が0でないので、この力には位置エネルギーが作れません。$\pxev F=(y,\ x)$ に変えて動かすと、0になります。

\begin{seikaku}
摩擦力は、いつも動く向きと反対にはたらくので、どんな道を通っても負の仕事をし、道が長いほど大きく負になります。だから、摩擦力は保存力ではありません。③で、摩擦の分が水平の距離 $\Delta x$ に比例して $v^2$ を減らしたのは、垂直抗力を $mg\cos\theta$ とした場合の特別な結果です。
\end{seikaku}

%% ---------------------------------------------------------------
\bksubsection{7}{力学的エネルギーの保存と、摩擦で失われる分}
\label{ph:力-4:sec7}
物にはたらく力を、保存力 $\pxev F_{\mathrm c}$ と、それ以外の力（摩擦力・空気の抵抗・人が押す力など）$\pxev F'$ に分けます。⑤から、
\[K_2-K_1=\int_C\pxev F_{\mathrm c}\cdot\dd\pxev r+\int_C\pxev F'\cdot\dd\pxev r=\bigl(U_1-U_2\bigr)+W'\]
です（$W'$ は保存力以外の力がした仕事）。$U$ を左へ移すと、
\pheq{(K_2+U_2)-(K_1+U_1)=W'}{J ＝ J}
\label{ph:力-4:sonshitsu}運動エネルギーと位置エネルギーの和 $E=K+U$ を\textbf{力学的エネルギー}といいます。\bkidx{りきがくてきえねるぎー}{力学的エネルギー}「力学的エネルギーの変化は、保存力以外の力がした仕事に等しい」。とくに、保存力以外の力が仕事をしないとき（$W'=0$）は、
\pheq{K+U=\frac{1}{2}mv^2+mgh=\text{一定}}{J}
\label{ph:力-4:hozon}です。これが\textbf{力学的エネルギー保存の法則}です。

\begin{phyomi}[$\dfrac{1}{2}mv^2+mgh=$ 一定]
「高さ（位置エネルギー）が減った分だけ、速さ（運動エネルギー）がふえる。高さがふえた分だけ、速さが減る」。①の表で $v^2$ が下りた高さに比例したのは、このため。
\end{phyomi}

\begin{bkfigure}
\begin{tikzpicture}
\begin{axis}[bkplot, width=118mm, height=58mm, xmin=0, xmax=108, ymin=0, ymax=470, xtick={20,40,60,80,100}, ytick={100,200,300,392}, xlabel={$x$（m）}, ylabel={エネルギー（J/kg）}]
  \fill[fill=bkOrange!25] (axis cs:0,392) -- (axis cs:10,333.2) -- (axis cs:20,215.6) -- (axis cs:30,98) -- (axis cs:40,29.4) -- (axis cs:50,19.6) -- (axis cs:60,78.4) -- (axis cs:70,176.4) -- (axis cs:80,274.4) -- (axis cs:90,343) -- (axis cs:100,372.4) -- (axis cs:100,0) -- (axis cs:0,0) -- cycle;
  \fill[fill=bkBlue!18] (axis cs:0,392) -- (axis cs:10,333.2) -- (axis cs:20,215.6) -- (axis cs:30,98) -- (axis cs:40,29.4) -- (axis cs:50,19.6) -- (axis cs:60,78.4) -- (axis cs:70,176.4) -- (axis cs:80,274.4) -- (axis cs:90,343) -- (axis cs:100,372.4) -- (axis cs:100,392) -- (axis cs:0,392) -- cycle;
  \draw[bkInk, line width=1.1pt] plot coordinates {(axis cs:0,392) (axis cs:10,333.2) (axis cs:20,215.6) (axis cs:30,98) (axis cs:40,29.4) (axis cs:50,19.6) (axis cs:60,78.4) (axis cs:70,176.4) (axis cs:80,274.4) (axis cs:90,343) (axis cs:100,372.4)};
  \draw[bkRed, line width=1.3pt] plot coordinates {(axis cs:0,392) (axis cs:100,392)};
  \draw[bkGreen!60!black, line width=1.3pt, dashed] plot coordinates {(axis cs:0,392) (axis cs:90,365.5) (axis cs:97,363.4)};
  \node[font=\sffamily\footnotesize, anchor=center] at (axis cs:85,110) {位置エネルギー $gh$};
  \node[font=\sffamily\footnotesize, anchor=center] at (axis cs:50,200) {運動エネルギー $\frac{1}{2}v^2$};
  \node[font=\sffamily\footnotesize, text=bkRed!75!black, anchor=south] at (axis cs:50,395) {摩擦なし：和は一定（392）};
  \node[font=\sffamily\footnotesize, text=bkGreen!50!black, anchor=north west] at (axis cs:24,368) {摩擦あり：和が減る（点線）};
\end{axis}
\end{tikzpicture}
\bkfigcaption{③のコースの、質量 \qty{1}{kg} あたりのエネルギー。だいだいの面が位置エネルギー、青い面が運動エネルギー。摩擦なしでは和がいつも $9.8\times40=392$。}\label{ph:力-4:fig-energy}
\end{bkfigure}

\phboxhead{\bkhatenaref の答え（式で）}
コースターにはたらく力は、重力（保存力）、垂直抗力（動く向きと直角なので仕事をしない）、摩擦力と空気の抵抗（いつも負の仕事 $W'<0$）です。だから、$K+U$ は増えることがなく、少しずつ減っていきます。最初の山の上（速さほぼ0）の $K+U=mgH$ より大きくなれないので、ある所で $mgh\leqq K+U<mgH$、つまり、\textbf{どの山も最初の山の高さ $H$ より低くしか上がれません}。最初の山より高く上がるには、モーターやチェーン（$W'>0$ の仕事をする力）で引き上げる必要があります。

\begin{phmijika}[ブレーキの熱]
摩擦力がした負の仕事の分の力学的エネルギーは、なくなったのではなく、熱（ブレーキやレールの温度が上がる）や音になっています。熱も入れたエネルギーの和は保存します。熱まで入れたエネルギーの数え方（熱力学の第一法則）は、熱-2 で学びます。
\end{phmijika}

\begin{phflow}[力学的エネルギー]
\phflowrow{$K_2-K_1=W$（すべての力）}{保存力の分を\\$U$ で書く}{$E_2-E_1=W'$}
\phflowrow{$E_2-E_1=W'$}{保存力以外は\\仕事をしない}{$\dfrac{1}{2}mv^2+mgh=$ 一定}
\end{phflow}

\begin{motto}[4]
\phboxhead{ばねにつけたおもりのエネルギー（波-1 の予告）}
なめらかな床の上で、ばね定数 $k$ のばねにつけたおもりを、自然の長さから $A$ だけ引いてはなすと、おもりは $-A$ と $A$ のあいだを行ったり来たりします（単振動。くわしくは波-1）。ばねの力は保存力なので、
\[\frac{1}{2}mv^2+\frac{1}{2}kx^2=\frac{1}{2}kA^2\qquad\text{（一定）}\]
です。たとえば、ばね定数が\qty{200}{N/m}で、\qty{0.10}{m}引いてはなしたときは、
\[\frac{1}{2}\times200\times0.10^2=\qty{1.0}{J}\]
が、いつも変わらない和です。はしでは全部が位置エネルギーに、まん中では全部が運動エネルギーになります。
\par\smallskip\noindent\begin{center}
\begin{tikzpicture}
\begin{axis}[bkplot, width=96mm, height=44mm, xmin=-0.12, xmax=0.12, ymin=0, ymax=1.3, xtick={-0.1,0,0.1}, xticklabels={$-A$,0,$A$}, ytick={0.5,1}, xlabel={$x$（m）}, ylabel={エネルギー（J）}, axis y line=left]
  \addplot[bkcurve, domain=-0.1:0.1] {100*x^2};
  \addplot[bkcurve2, domain=-0.1:0.1] {1-100*x^2};
  \draw[bkRed, line width=1.2pt] (axis cs:-0.1,1) -- (axis cs:0.1,1);
  \node[font=\sffamily\footnotesize, text=bkBlue!80!black, anchor=south west, fill=white, inner sep=1pt] at (axis cs:0.084,0.4) {$\frac{1}{2}kx^2$};
  \node[font=\sffamily\footnotesize, text=bkOrange!75!black, anchor=west, fill=white, inner sep=1pt] at (axis cs:-0.035,0.7) {$\frac{1}{2}mv^2$};
  \node[font=\sffamily\footnotesize, text=bkRed!75!black, anchor=south] at (axis cs:0,1.02) {和は一定（1.0 J）};
\end{axis}
\end{tikzpicture}
\end{center}
\end{motto}

\begin{phsaki}{エネルギーは、なぜ保存するのか}
1918年、数学者のエミー・ネーターは、「自然の法則が、時間がたっても変わらない（きのうの実験と今日の実験で、同じ法則が成り立つ）」ことから、エネルギーの保存が導けることを証明しました（ネーターの定理）。同じように、「場所を移しても法則が変わらない」ことから運動量の保存が、「向きを変えても法則が変わらない」ことから角運動量の保存（力-6）が導かれます。保存則と、自然の「対称性」が結びついているのです。くわしくは、大学の解析力学で学びます。
\end{phsaki}

%% ---------------------------------------------------------------
\bksubsection{8}{仕事率・動力・効率・仕事の原理}
\label{ph:力-4:sec8}
\phboxhead{仕事率：1秒あたりの仕事}
同じ荷物を同じ高さに上げても、すばやく上げるほうが、強いモーターがいります。1秒あたりにする仕事を\textbf{仕事率}（機械では\textbf{動力}）といい、$P$ で表します。単位は W（ワット）で、\qty{1}{W}＝\qty{1}{J/s}です。\bkidx{しごとりつ}{仕事率}力 $\pxev F$ が、少しの時間 $\dd t$ に $\dd\pxev r$ 動かすときの仕事は $\pxev F\cdot\dd\pxev r$ なので、
\pheq{P=\frac{\dd W}{\dd t}=\pxev F\cdot\frac{\dd\pxev r}{\dd t}=\pxev F\cdot\pxev v}{N × m/s ＝ J/s ＝ W}
\label{ph:力-4:power}です。仕事率が一定なら $P=\dfrac{W}{t}$、力と速度が同じ向きなら $P=Fv$ です。

\begin{phbaisuru}{速さ $v$（力は同じ）}
同じ力で2倍の速さで動かすと、仕事率は2倍です。坂道を同じ自転車で2倍の速さで上ると、坂を上るための仕事率も2倍いります（例題8）。
\end{phbaisuru}

\phboxhead{効率：入れたエネルギーのうち、役に立った割合}
\label{ph:力-4:kouritsu}ポンプやモーターは、入れたエネルギー（電気など）の一部を、摩擦や熱でむだにします。役に立った仕事率 $P_{\text{出}}$ と、入れた仕事率 $P_{\text{入}}$ の比を\textbf{効率} $\eta$（エータ）といいます。
\pheq{\eta=\frac{P_{\text{出}}}{P_{\text{入}}}}{W÷W：単位なし}
$\eta$ は1（100\,\%）より小さい数です。たとえば $\eta=0.60$ のポンプに\qty{100}{W}を入れると、水を上げる仕事に使われるのは\qty{60}{W}で、残りの\qty{40}{W}は熱や音になります。

\phboxhead{仕事の原理：道具を使っても、仕事は得をしない}
\label{ph:力-4:genri}\noindent\begin{minipage}[c]{0.46\linewidth}\centering
\begin{tikzpicture}[x=0.82mm, y=0.82mm]
  \fill[bkGray!25] (0,0) -- (48,0) -- (48,24) -- cycle;
  \draw[bkInk!70, line width=0.7pt] (0,0) -- (48,0) -- (48,24) -- cycle;
  \pxebox{(20,10)}{7}{6}{bkOrange}
  \draw[bk arrow, bkBlue] (13,6.5) -- (21,10.5);
  \node[font=\sffamily\footnotesize, anchor=north west, text=bkBlue!80!black] at (13,5) {小さい力・長い距離};
  \draw[bk arrow, bkRed] (54,0) -- (54,22);
  \node[font=\sffamily\footnotesize, anchor=west, text=bkRed!75!black, align=left] at (55,11) {大きい力\\短い距離};
  \draw[bkInk, line width=0.5pt] (40,0) arc (180:153.4:8);
  \node[font=\sffamily\footnotesize] at (37.5,2.4) {$\theta$};
\end{tikzpicture}
\end{minipage}\hfill
\begin{minipage}[c]{0.51\linewidth}\raggedright
重さ $mg$ の荷物を高さ $h$ まで上げます。まっすぐ持ち上げると、力 $mg$ で距離 $h$。なめらかな斜面を使うと、力 $mg\sin\theta$ で距離 $\dfrac{h}{\sin\theta}$。どちらも仕事は $mgh$ です。
\end{minipage}

\par\vspace{4pt}
摩擦も、道具そのものの質量もない理想の道具では、道具に入れた仕事と、道具が荷物にした仕事は等しくなります（道具はエネルギーをためも、むだにもしないから）。
\pheq{F_1s_1=F_2s_2}{J ＝ J}
斜面・てこ・滑車・ねじは、力を小さくできますが、そのかわりに、動かす距離が長くなります。これを\textbf{仕事の原理}といいます。本物の道具には摩擦があるので、入れる仕事のほうが大きくなります（効率が1より小さい）。滑車やてこを、力のつり合い（モーメント）で調べる方法は、力-5 で学びます。

\begin{phflow}[仕事率と道具]
\phflowrow{$P=\pxev F\cdot\pxev v$}{力と速度が\\同じ向き}{$P=Fv$}
\phflowrow{$P=\dfrac{\dd W}{\dd t}$}{仕事率が\\一定}{$P=\dfrac{W}{t}$}
\phflowrow{入れた仕事＝\\した仕事＋失った分}{摩擦・道具の\\質量がない}{$F_1s_1=F_2s_2$}
\end{phflow}

%% ---------------------------------------------------------------
\bksubsection{9}{ループを1周する：鉛直面内の円運動}
\label{ph:力-4:sec9}
\noindent\begin{minipage}[c]{0.44\linewidth}\centering
\begin{tikzpicture}[x=0.82mm, y=0.82mm]
  \draw[bkInk, line width=1.1pt] (14,14) circle (12);
  \draw[bkInk, line width=1.1pt] plot[smooth, tension=0.5] coordinates {(-12,40) (-6,34) (2,8) (8,2.2) (14,2)};
  \draw[bkInk, line width=1.1pt] (14,2) -- (40,2);
  \draw[bkGray, line width=0.4pt, dashed] (-14,40) -- (44,40) (-14,2) -- (44,2);
  \draw[<->, bkBlue, line width=0.6pt] (42,2.5) -- (42,39.5);
  \node[font=\sffamily\footnotesize, anchor=west, text=bkBlue!80!black] at (42.5,21) {$H$};
  \draw[bkInk, line width=0.5pt] (14,14) -- (26,14);
  \fill[bkInk] (14,14) circle (0.6);
  \node[font=\sffamily\footnotesize, anchor=north] at (20,13.5) {$r$};
  \fill[bkRed] (14,26) circle (1.2);
  \draw[bk arrow, bkPurple, line width=1pt] (14,24.6) -- (14,17.5);
  \node[font=\sffamily\footnotesize, anchor=west, text=bkPurple!80!black] at (15,21) {$mg+N$};
  \node[font=\sffamily\footnotesize, anchor=south, text=bkRed!75!black] at (14,27.5) {いちばん上};
  \fill[bkRed] (-12,40) circle (1.2);
\end{tikzpicture}
\end{minipage}\hfill
\begin{minipage}[c]{0.53\linewidth}\raggedright
高さ $H$ から速さ0で下り始めた車が、半径 $r$ のループを、落ちずに1周するには、$H$ がどれだけいるでしょう（摩擦なし）。\par\vspace{2pt}
いちばん上では、レールは車の上にあり、レールが車を押す力 $N$（下向き）と重力 $mg$（下向き）が、どちらも円の中心向きです。
\end{minipage}

\par\vspace{4pt}
いちばん上で、中心向きの運動方程式（力-2の円運動）は
\[m\frac{v^2}{r}=mg+N\]
です。レールから離れないためには $N\geqq0$、つまり $v^2\geqq gr$ です。いちばん上の高さは $2r$ なので、力学的エネルギー保存から、
\[mgH=\frac{1}{2}mv^2+mg\cdot2r\qquad v^2=2g(H-2r)\]
$v^2\geqq gr$ に入れると、$2g(H-2r)\geqq gr$、つまり
\pheq{H\geqq\frac{5}{2}r}{m ＝ m}
\label{ph:力-4:loop}です。ループの半径の2.5倍の高さから下りれば、1周できます。いちばん下の速さで言えば、$v^2=2gH\geqq5gr$ です。

\begin{phyomi}[$H\geqq\dfrac{5}{2}r$]
「ループのいちばん上の高さ $2r$ に、そこで落ちないための速さの分 $\dfrac{r}{2}$（$\dfrac{v^2}{2g}=\dfrac{gr}{2g}$）をたした高さがいる」。
\end{phyomi}

\begin{phbaisuru}{ループの半径 $r$}
いる高さ $H$ も2倍になります。そのうえ、いちばん下でレールが車を押す力は、$H=2.5r$ のとき $N=m\Bigl(g+\dfrac{v^2}{r}\Bigr)=m(g+5g)=6mg$ で、$r$ によりません。乗る人は、体重の6倍の力で座席に押しつけられます。そこで本物のループは、下を大きく、上を小さくした「しずく形」にして、この力を小さくしています。
\end{phbaisuru}

\begin{phmada}[転がる・こすれるときに、エネルギーはどう熱になる？]
車輪が転がると、タイヤやレールはほんの少しへこみ、またもとにもどります。そのたびにエネルギーの一部が熱になります（転がり抵抗）。空気の抵抗も、物のまわりの空気の流れでエネルギーが熱になる現象です。このような抵抗の大きさは、ゴムの性質や空気の流れから理論だけで正確に予測することは今もむずかしく、ふつうは実験で係数を決めて使っています（空気の抵抗は力-15）。摩擦（力-2）と同じく、「エネルギーが熱に変わる所」のくわしいしくみは、今も研究が続いています。
\end{phmada}

\begin{senmonsho}[力学の専門書の書き方]
\begin{bkhonbun}
力 $\symbf F$ が保存力であるための必要十分条件は，$\nabla\times\symbf F=\symbf 0$ が単連結な領域で成り立つことである．このときポテンシャル $U$ が存在して $\symbf F=-\nabla U$ と書け，力学的エネルギー $E=\dfrac{1}{2}m|\dot{\symbf r}|^2+U(\symbf r)$ は保存される．
\end{bkhonbun}
\bkyaku{太字の $\symbf F$ は、この本の矢印 $\vec F$ と同じベクトル（4-13）。$\nabla\times\symbf F$ は $\mathrm{rot}\,\pxev F$、$\nabla U$ は $\mathrm{grad}\,U$ の書き方（8-4・8-10）。「単連結な領域」は、穴のない広がり（⑥の「すきまのない広がり」）。$\dot{\symbf r}$ は速度。この章の⑥と⑦を1つの文にしたもの。}
\end{senmonsho}

\begin{taisetsu}
\def\tsl#1{\makebox[6.4\zw][l]{\sffamily\bfseries\small\color{\bkTT}#1}}
\tsl{仕事}$W=\displaystyle\int_C\pxev F\cdot\dd\pxev r$。一定なら $Fs\cos\theta$\\[3pt]
\tsl{運動エネルギー}$K=\dfrac{1}{2}mv^2$。$K_2-K_1=W$\\[3pt]
\tsl{位置エネルギー}$\pxev F=-\mathrm{grad}\,U$。$mgh$、$\dfrac{1}{2}kx^2$、$-\dfrac{GMm}{r}$\\[3pt]
\tsl{保存の法則}$K+U=$ 一定（保存力だけが仕事をするとき）\\[3pt]
\tsl{失われる分}$E_2-E_1=W'$（摩擦などの仕事）\\[3pt]
\tsl{仕事率}$P=\pxev F\cdot\pxev v$。効率 $\eta=\dfrac{P_{\text{出}}}{P_{\text{入}}}$
\end{taisetsu}

\begin{kotoba}
\small
\bkword{仕事}{力×力の向きに動いた距離。一般には線積分。単位 J。}
\bkword{運動エネルギー}{$\dfrac{1}{2}mv^2$。された仕事の分だけ変わる。}
\bkword{保存力}{仕事が、はじめとおわりの点だけで決まる力（重力・ばね・万有引力）。}
\bkword{位置エネルギー}{保存力について、点ごとに決まるエネルギー。力はそれが下がる向き。}
\bkword{力学的エネルギー}{運動エネルギーと位置エネルギーの和。}
\bkword{仕事率（動力）}{1秒あたりの仕事。単位 W＝J/s。}
\end{kotoba}

\bkkigocards{\bkkigocard{$U$}{ユー}{位置エネルギー。力は $-\mathrm{grad}\,U$。$K+U=E$（力学的エネルギー）}{potential energy}{力-4}}
{\bkkigocard{$P$}{ピー}{仕事率・動力（W）。荷重（力-11）や電力にも使う}{power}{力-4}}

%% ======================== 例題 ========================
\Needspace{10\baselineskip}\bksection{reidai}{例題でたしかめよう}
\label{ph:力-4:ex1}
\begin{phreidai}[ジェットコースターの谷での速さ]
高さ\qty{60}{m}の最初の山の上から、ほぼ速さ0で下り始めたジェットコースターが、いちばん低い谷（高さ0）を通るときの速さを求めましょう。摩擦と空気の抵抗は無視します。
\phbamen
高さが速さに変わる動きです。コースの形はわかりませんが、高さはわかっています。
\phmushi
摩擦・空気の抵抗・車輪の回転のエネルギーを無視します。車は点と考えます。
\phzu
谷を高さの基準（0）にします。山の上：高さ $H=\qty{60}{m}$、速さ0。谷：高さ0、速さ $v$。
\phriyu
重力は保存力で、垂直抗力は動く向きと直角なので仕事をしません。保存力以外の力が仕事をしないので、力学的エネルギー保存（⑦）が使えます。コースの形を知らなくても、高さだけで答えが出ます。
\phtateru
\[0+mgH=\frac{1}{2}mv^2+0\]
\phtoku
両辺を $m$ でわって、
\[v=\sqrt{2gH}=\sqrt{2\times9.8\times60}=\sqrt{1176}\approx\qty{34}{m/s}\]
\phtashikame
単位：$\sqrt{\text{m/s}^2\times\text{m}}=$ m/s。\\
大きさ：\qty{34}{m/s}は時速約\qty{123}{km}です。本物の大きなコースターの最高速度と同じくらいの桁ですが、本物は摩擦と空気の抵抗の分だけおそくなります。\\
質量：$m$ が消えたので、乗っている人数によらず、同じ速さです（摩擦がなければ）。
\kotae 約\qty{34}{m/s}
\end{phreidai}
\begin{jibun}
高さ\qty{50}{m}の山の上から速さ0で下り始めたコースターが、高さ\qty{20}{m}の所を通るときの速さを求めましょう。摩擦と空気の抵抗は無視します。
\bkjans{$\sqrt{2\times9.8\times30}\approx\qty{24}{m/s}$}
\end{jibun}

\label{ph:力-4:ex2}
\begin{phreidai}[ループを1周するのに必要な高さ]
半径\qty{10}{m}の円のループを、コースターが落ちずに1周するには、速さ0で下り始める山の高さは少なくとも何mいりますか。また、そのとき、いちばん下で座席が人（質量 $m$）を押す力は、体重の何倍ですか。
\phbamen
ループのいちばん上では、車はさかさまになります。速さが足りないと、レールから離れて落ちてしまいます。
\phmushi
摩擦と空気の抵抗を無視し、車を点と考えます。ループは完全な円とします。
\phzu
⑨の図と同じです。山の高さ $H$、ループの半径 $r=\qty{10}{m}$、いちばん上での速さ $v$、レールが押す力 $N$。
\phriyu
「落ちない条件」は力の条件なので、いちばん上で円運動の運動方程式（力-2）を使います。速さは高さで決まるので、力学的エネルギー保存（⑦）を使います。2つを組み合わせます。
\phtateru
\[m\frac{v^2}{r}=mg+N,\quad N\geqq0\qquad\qquad mgH=\frac{1}{2}mv^2+mg\cdot2r\]
\phtoku
1つめの式から $v^2\geqq gr$、2つめの式から $v^2=2g(H-2r)$ です。だから、
\[2g(H-2r)\geqq gr\qquad H\geqq2.5r=2.5\times10=\qty{25}{m}\]
$H=\qty{25}{m}$ のとき、いちばん下の速さは $v_0^2=2gH=2\times9.8\times25=490$。座席が押す力は、中心向き（上向き）の運動方程式 $m\dfrac{v_0^2}{r}=N_0-mg$ から、
\[N_0=m\Bigl(g+\frac{490}{10}\Bigr)=m(9.8+49)=m\times58.8=6\times mg\]
\phtashikame
大きさ：ループのいちばん上の高さ（\qty{20}{m}）より、\qty{5}{m}だけ高い所から下りればよい。\\
極端な場合：$H$ がちょうど\qty{25}{m}なら、いちばん上で $N=0$、人は一瞬「重さがない」状態になります。
\kotae 少なくとも\qty{25}{m}。いちばん下で体重の6倍
\end{phreidai}
\begin{jibun}
半径\qty{8.0}{m}のループを、摩擦なしで1周するには、速さ0で下り始める高さは少なくとも何mいりますか。
\bkjans{$2.5\times8.0=\qty{20}{m}$}
\end{jibun}

\label{ph:力-4:ex3}
\begin{phreidai}[荷物を斜面で押し上げる]
質量\qty{20}{kg}の荷物を、角 $30^\circ$ のなめらかな斜面にそって、ゆっくり\qty{4.0}{m}押し上げます。斜面にそって押す力の大きさと、した仕事を求めましょう。また、同じ荷物を真上に持ち上げるときとくらべましょう。
\phbamen
斜面という道具を使って、荷物を高い所へ上げます。
\phmushi
斜面の摩擦を無視します。「ゆっくり」なので、運動エネルギーの変化は0とします。
\phzu
⑧の斜面の図と同じです。$m=\qty{20}{kg}$、$\theta=30^\circ$、押す距離 $s=\qty{4.0}{m}$、押す力 $F$、上がった高さ $h=s\sin\theta$。
\phriyu
ゆっくりなので、斜面にそった力はつり合っています（$F=mg\sin\theta$）。仕事は力×距離（④）、比べるのは仕事の原理（⑧）です。
\phtateru
\[F=mg\sin30^\circ,\qquad W=Fs\]
\phtoku
\[F=20\times9.8\times0.5=\qty{98}{N},\qquad W=98\times4.0=\qty{392}{J}\]
真上に持ち上げると、力 $20\times9.8=\qty{196}{N}$、高さ $h=4.0\times0.5=\qty{2.0}{m}$ で、$W=196\times2.0=\qty{392}{J}$。
\phtashikame
仕事の原理：力は半分、距離は2倍で、仕事は同じ\qty{392}{J}です。\\
エネルギーで：位置エネルギーのふえた分 $mgh=20\times9.8\times2.0=\qty{392}{J}$ と同じです。\\
摩擦があると：押す力は大きくなり、仕事は\qty{392}{J}より多くいります（効率が1より小さい）。
\kotae 力\qty{98}{N}、仕事\qty{392}{J}（真上に持ち上げても仕事は同じ\qty{392}{J}）
\end{phreidai}
\begin{jibun}
質量\qty{15}{kg}の荷物を、角 $30^\circ$ のなめらかな斜面にそって、ゆっくり\qty{6.0}{m}押し上げます。押す力と仕事を求めましょう。
\bkjans{力 $15\times9.8\times0.5=\qty{73.5}{N}$、仕事 $73.5\times6.0=\qty{441}{J}$}
\end{jibun}

\label{ph:力-4:ex4}
\begin{phreidai}[おもちゃの鉄砲]
ばね定数\qty{200}{N/m}のばねを\qty{5.0}{cm}縮めて、質量\qty{10}{g}の玉をのせ、ばねをはなしました。玉が飛び出す速さを求めましょう。
\phbamen
ばねにたくわえたエネルギーが、玉の運動エネルギーに変わります。
\phmushi
ばねの質量・摩擦・空気の抵抗を無視します。玉は水平に飛び出すとします（高さは変わらない）。
\phzu
$k=\qty{200}{N/m}$、縮み $x=\qty{5.0}{cm}=\qty{0.050}{m}$、$m=\qty{10}{g}=\qty{0.010}{kg}$、飛び出す速さ $v$。
\phriyu
ばねの力は保存力なので、ばねの位置エネルギー $\dfrac{1}{2}kx^2$（⑥）と運動エネルギーの和が保存します（⑦）。
\phtateru
\[\frac{1}{2}kx^2=\frac{1}{2}mv^2\]
\phtoku
\[v=x\sqrt{\frac{k}{m}}=0.050\times\sqrt{\frac{200}{0.010}}=0.050\times141.4\approx\qty{7.1}{m/s}\]
\phtashikame
単位：$\text{m}\times\sqrt{\dfrac{\text{N/m}}{\text{kg}}}=\text{m}\times\sqrt{\dfrac{1}{\text{s}^2}}=$ m/s。\\
大きさ：ばねのエネルギーは $\dfrac{1}{2}\times200\times0.050^2=\qty{0.25}{J}$ で、小さなおもちゃらしい値です。\\
2倍にすると：縮みを2倍にすると、速さも2倍（$v$ は $x$ に比例）。
\kotae 約\qty{7.1}{m/s}
\end{phreidai}
\begin{jibun}
ばね定数\qty{400}{N/m}のばねを\qty{4.0}{cm}縮め、質量\qty{20}{g}の玉を飛ばします。飛び出す速さを求めましょう。
\bkjans{$0.040\times\sqrt{400\div0.020}\approx\qty{5.7}{m/s}$}
\end{jibun}

\label{ph:力-4:ex5}
\begin{phreidai}[タイヤがすべって止まるまでの距離]
速さ\qty{20}{m/s}で走る車が、急ブレーキでタイヤがすべり（ロックし）、止まりました。タイヤと路面の動摩擦係数を0.70として、すべった距離を求めましょう。
\phbamen
動摩擦力が、車の運動エネルギーを熱に変えて、車を止めます。
\phmushi
空気の抵抗を無視し、路面は水平とします。動摩擦係数は速さによらず一定とします。
\phzu
$v=\qty{20}{m/s}$、$\mu'=0.70$、すべった距離 $d$。動摩擦力の大きさ $\mu'mg$。
\phriyu
動摩擦力は保存力ではないので、「力学的エネルギーの変化＝保存力以外の力の仕事」（⑦）を使います。動摩擦力は動く向きと反対向きなので、仕事は $-\mu'mg\,d$ です。
\phtateru
\[0-\frac{1}{2}mv^2=-\mu'mg\,d\]
\phtoku
\[d=\frac{v^2}{2\mu'g}=\frac{20^2}{2\times0.70\times9.8}=\frac{400}{13.72}\approx\qty{29}{m}\]
\phtashikame
単位：$\dfrac{\text{m}^2/\text{s}^2}{\text{m/s}^2}=$ m。\\
大きさ：力-1の、減速 $\mu'g\approx\qty{6.9}{m/s^2}$ で止まる距離 $\dfrac{v^2}{2a}$ と同じです。\\
2倍にすると：速さが2倍（\qty{40}{m/s}）なら、距離は4倍の約\qty{117}{m}です。$m$ が消えたので、重い車でも軽い車でも、同じ $\mu'$ なら同じ距離です。
\kotae 約\qty{29}{m}
\end{phreidai}
\begin{jibun}
速さ\qty{25}{m/s}の車が急ブレーキですべって止まりました。動摩擦係数0.70として、すべった距離を求めましょう。
\bkjans{$25^2\div(2\times0.70\times9.8)\approx\qty{46}{m}$}
\end{jibun}

\label{ph:力-4:ex6}
\begin{phreidai}[クレーンのモーターの動力]
クレーンが、質量\qty{2.0}{t}の荷物を、一定の速さ\qty{0.50}{m/s}で真上に引き上げています。ワイヤが荷物を引く仕事率を求めましょう。
\phbamen
一定の速さで上げるので、ワイヤの張力は重さとつり合っています。
\phmushi
ワイヤの質量、滑車の摩擦、モーターの中のむだ（効率）を無視します。
\phzu
$m=\qty{2.0e3}{kg}$、$v=\qty{0.50}{m/s}$（上向き）、張力 $T$（上向き）、仕事率 $P$。
\phriyu
力と速度が同じ向きなので、⑧の $P=Fv$ が使えます。
\phtateru
\[T=mg,\qquad P=Tv\]
\phtoku
\[P=2000\times9.8\times0.50=\qty{9.8e3}{W}=\qty{9.8}{kW}\]
\phtashikame
別の見方：1秒に\qty{0.50}{m}上がるので、1秒あたり位置エネルギーが $2000\times9.8\times0.50=\qty{9800}{J}$ ふえます。\\
大きさ：\qty{9.8}{kW}は、家庭のエアコン数台分の電力の桁です。本物のモーターは効率の分だけ大きい動力が要ります。
\kotae \qty{9.8}{kW}
\end{phreidai}
\begin{jibun}
エレベータが、質量\qty{600}{kg}のかごを一定の速さ\qty{1.5}{m/s}で引き上げています。ロープが引く仕事率を求めましょう。
\bkjans{$600\times9.8\times1.5=\qty{8.82e3}{W}$（約\qty{8.8}{kW}）}
\end{jibun}

\label{ph:力-4:ex7}
\begin{phreidai}[階段をかけ上がる人の仕事率]
質量\qty{60}{kg}の人が、高さ\qty{3.0}{m}の階段を\qty{5.0}{s}でかけ上がりました。この人が体を持ち上げた仕事率を求めましょう。
\phbamen
人の体が「モーター」になって、自分を持ち上げます。
\phmushi
階段の上と下で速さは同じとし、運動エネルギーの変化は考えません。体の中で熱になる分も考えません（体を持ち上げる仕事だけ）。
\phzu
$m=\qty{60}{kg}$、$h=\qty{3.0}{m}$、$t=\qty{5.0}{s}$。
\phriyu
した仕事は位置エネルギーのふえた分 $mgh$ です（⑥）。かかった時間でわれば、平均の仕事率です（⑧）。
\phtateru
\[P=\frac{W}{t}=\frac{mgh}{t}\]
\phtoku
\[P=\frac{60\times9.8\times3.0}{5.0}=\frac{1764}{5.0}\approx\qty{3.5e2}{W}\]
\phtashikame
大きさ：人が8時間の仕事で続けられる仕事率は約\qty{75}{W}です（工学ではこう使う）。\qty{350}{W}は、短い時間だけ出せる値です。\\
2倍にすると：同じ階段を半分の時間（\qty{2.5}{s}）で上ると、仕事率は2倍です。
\kotae 約\qty{3.5e2}{W}
\end{phreidai}
\begin{jibun}
質量\qty{50}{kg}の人が、高さ\qty{4.0}{m}の階段を\qty{8.0}{s}で上りました。体を持ち上げた仕事率を求めましょう。
\bkjans{$50\times9.8\times4.0\div8.0\approx\qty{2.5e2}{W}$}
\end{jibun}

\label{ph:力-4:ex8}
\begin{phreidai}[自転車で坂を上る]
人と自転車を合わせて質量\qty{70}{kg}が、こう配5\,\%の坂を、一定の速さ\qty{4.0}{m/s}で上っています。坂を上るためだけに必要な仕事率を求めましょう。
\phbamen
こう配5\,\%は、水平に\qty{100}{m}進むと\qty{5}{m}上がる坂です（4-6）。
\phmushi
空気の抵抗とタイヤの転がる抵抗を無視します（坂を上る分だけを求める）。
\phzu
$m=\qty{70}{kg}$、$\tan\theta=0.05$、$v=\qty{4.0}{m/s}$（坂にそって）。坂にそって下向きの重力の成分 $mg\sin\theta$。
\phriyu
一定の速さなので、こぐ力は $mg\sin\theta$ とつり合っています。力と速度が同じ向きなので $P=Fv$（⑧）。
\phtateru
\[P=mg\sin\theta\times v\]
\phtoku
$\tan\theta=0.05$ のとき、$\sin\theta=\dfrac{0.05}{\sqrt{1+0.05^2}}\approx0.0499$（角が小さいので $\sin\theta\approx\tan\theta$）。
\[P=70\times9.8\times0.0499\times4.0\approx\qty{1.4e2}{W}\]
\phtashikame
別の見方：1秒に上がる高さは $4.0\times0.0499\approx\qty{0.20}{m}$。1秒あたりの位置エネルギーのふえ方 $70\times9.8\times0.20\approx\qty{137}{J}$ と同じです。\\
大きさ：人が長く続けられる約\qty{75}{W}より大きいので、この坂を\qty{4.0}{m/s}で上り続けるのはたいへんです。
\kotae 約\qty{1.4e2}{W}
\end{phreidai}
\begin{jibun}
人と自転車で質量\qty{80}{kg}が、こう配8\,\%の坂を一定の速さ\qty{3.0}{m/s}で上ります。坂を上るためだけの仕事率を求めましょう（$\sin\theta\approx\tan\theta$ としてよい）。
\bkjans{$80\times9.8\times0.08\times3.0\approx\qty{1.9e2}{W}$}
\end{jibun}

\label{ph:力-4:ex9}
\begin{phreidai}[ポンプで水をくみ上げる]
ポンプで、水を毎分\qty{100}{L}ずつ、\qty{10}{m}の高さまでくみ上げます。水をくみ上げる仕事率（有効な動力）と、ポンプの効率を0.60としたときの、ポンプに入れる動力を求めましょう。
\phbamen
ポンプは、水に位置エネルギーを与える機械です。
\phmushi
管の中の摩擦による損失（力-14）と、出口の水の運動エネルギーを無視します。
\phzu
流量 $Q=\qty{100}{L/min}$、高さ $H=\qty{10}{m}$、水の密度 $\rho=\qty{1000}{kg/m^3}$、効率 $\eta=0.60$。
\phriyu
1秒あたりに上げる水の質量 $\rho Q$ に、$gH$ をかければ、1秒あたりにふえる位置エネルギー（有効な動力）です。入れる動力は、効率でわります（⑧）。
\phtateru
\[P_{\text{出}}=\rho QgH,\qquad P_{\text{入}}=\frac{P_{\text{出}}}{\eta}\]
\phtoku
$Q=\qty{100}{L/min}=\dfrac{\qty{0.100}{m^3}}{\qty{60}{s}}\approx\qty{1.667e-3}{m^3/s}$（\qty{1000}{L/m^3}、\qty{60}{s/min}で換算）。
\[P_{\text{出}}=1000\times\num{1.667e-3}\times9.8\times10\approx\qty{163}{W}\]
\[P_{\text{入}}=\frac{163}{0.60}\approx\qty{272}{W}\]
\phtashikame
効率の向き：入れる動力のほうが大きくなるのが正しい（$0.60$ をかけると小さくなってしまう）。\\
大きさ：家庭の井戸ポンプの電力（数百 W）と同じ桁です。
\kotae 有効な動力 約\qty{163}{W}、入れる動力 約\qty{272}{W}
\end{phreidai}
\begin{jibun}
毎分\qty{60}{L}の水を\qty{5.0}{m}の高さまでくみ上げます。有効な動力と、効率0.50のときに入れる動力を求めましょう。
\bkjans{有効な動力 $1000\times\num{1.0e-3}\times9.8\times5.0=\qty{49}{W}$、入れる動力\qty{98}{W}}
\end{jibun}

\label{ph:力-4:ex10}
\begin{phreidai}[水力発電所の出力]
水車に、毎秒\qty{10}{m^3}の水が、高さの差\qty{100}{m}の所から流れこみます。水車と発電機を合わせた効率を0.85として、発電所の出力を求めましょう。
\phbamen
ポンプの反対で、水の位置エネルギーを、水車で電気のエネルギーに変えます。
\phmushi
水路の摩擦による損失を、効率の中に入れて考えます。
\phzu
$Q=\qty{10}{m^3/s}$、高さの差 $H=\qty{100}{m}$、$\eta=0.85$。
\phriyu
1秒あたりに下りる水の位置エネルギー $\rho QgH$ に、効率をかけたものが、1秒あたりにできる電気のエネルギー（出力）です（⑧）。
\phtateru
\[P=\eta\,\rho QgH\]
\phtoku
\[P=0.85\times1000\times10\times9.8\times100=\qty{8.33e6}{W}\approx\qty{8.3}{MW}\]
\phtashikame
工学の式：$\rho g=1000\times9.8=9800$ なので、$P$（kW）$\approx9.8\times Q\times H\times\eta$ と書けます（$Q$ は m$^3$/s、$H$ は m）。$9.8\times10\times100\times0.85=\num{8330}$ kW で同じです。この形は力-13 でも使います。\\
こんどは効率をかける：エネルギーを取り出すので、出てくる電気は水のエネルギーより小さくなります。
\kotae 約\qty{8.3}{MW}
\end{phreidai}
\begin{jibun}
毎秒\qty{2.0}{m^3}の水が高さの差\qty{50}{m}を下りる発電所の出力を、効率0.80として求めましょう。
\bkjans{$0.80\times1000\times2.0\times9.8\times50$ より、約\qty{0.78}{MW}}
\end{jibun}

\label{ph:力-4:ex11}
\begin{phreidai}[地球から飛び出す速さ]
地面から真上に打ち出した物が、地球にもどってこないための、いちばん小さい速さ（第2宇宙速度）を、3けたで求めましょう。地球の $GM$ は\qty{3.986e14}{m^3/s^2}、半径 $R$ は\qty{6.371e6}{m}とします。
\phbamen
高く上がると、万有引力が弱くなるので、$mgh$ ではなく $-\dfrac{GMm}{r}$ を使います。
\phmushi
空気の抵抗と、地球の自転、ほかの天体（月・太陽）の引力を無視します。
\phzu
地球の中心からの距離 $r$。打ち出す所 $r=R$、速さ $v$。無限に遠い所で速さ0。
\phriyu
万有引力は保存力なので、力学的エネルギー保存（⑦）と、万有引力の位置エネルギー（⑥）を使います。
\phtateru
\[\frac{1}{2}mv^2-\frac{GMm}{R}=0\]
\phtoku
\[v=\sqrt{\frac{2GM}{R}}=\sqrt{\frac{2\times\num{3.986e14}}{\num{6.371e6}}}=\sqrt{\num{1.251e8}}\approx\qty{1.12e4}{m/s}\]
\phtashikame
単位：$\sqrt{\dfrac{\text{m}^3/\text{s}^2}{\text{m}}}=$ m/s。\\
大きさ：\qty{11.2}{km/s}は、人工衛星の速さ約\qty{7.9}{km/s}（力-2）の $\sqrt2$ 倍です。\\
$m$ が消えた：重い物も軽い物も、同じ速さがいります（でも、重い物ほど、その速さにするエネルギーは大きい）。
\kotae 約\qty{11.2}{km/s}
\end{phreidai}
\begin{jibun}
月の $GM=\qty{4.905e12}{m^3/s^2}$、半径 $\qty{1.737e6}{m}$ として、月の表面から飛び出す速さを求めましょう。
\bkjans{$v^2=2\times\num{4.905e12}\div\num{1.737e6}$ より、約\qty{2.4}{km/s}}
\end{jibun}

\label{ph:力-4:ex12}
\begin{phreidai}[同じ高さを2つの道で下りる]
質量\qty{10}{kg}のそりが、高さ\qty{5.0}{m}の所から速さ0で、（A）角 $30^\circ$ の坂と、（B）角 $10^\circ$ の坂を、それぞれすべり下ります。動摩擦係数を0.10として、下に着いたときの速さを求めましょう。
\phbamen
始まりと終わりの高さは同じで、道だけがちがいます。
\phmushi
空気の抵抗を無視し、坂はまっすぐとします。
\phzu
③の図と同じです。坂の長さ $s=\dfrac{h}{\sin\theta}$、水平の距離 $\Delta x=\dfrac{h}{\tan\theta}$。
\phriyu
重力（保存力）の仕事は道によらず $mgh$。摩擦力（保存力でない）の仕事は道で変わり、$-\mu'mg\cos\theta\times s=-\mu'mg\,\Delta x$（③）。⑤の「運動エネルギーの変化＝仕事の和」を使います。
\phtateru
\[\frac{1}{2}mv^2=mgh-\mu'mg\,\Delta x\]
\phtoku
（A）$\Delta x=\dfrac{5.0}{\tan30^\circ}\approx\qty{8.66}{m}$、摩擦の仕事 $-0.10\times10\times9.8\times8.66\approx\qty{-85}{J}$、$\dfrac{1}{2}\times10v^2=490-85=405$、$v\approx\qty{9.0}{m/s}$。\\
（B）$\Delta x=\dfrac{5.0}{\tan10^\circ}\approx\qty{28.4}{m}$、摩擦の仕事 $\approx\qty{-278}{J}$、$5v^2=490-278=212$、$v\approx\qty{6.5}{m/s}$。
\phtashikame
道の比べ方：重力の仕事は両方\qty{490}{J}で同じ、摩擦の仕事は長い道（B）のほうが大きく負です。仕事が道で変わる力（摩擦）と、変わらない力（重力）のちがいです。\\
極端な場合：$\mu'=0$ なら、どちらも $\sqrt{2\times9.8\times5.0}\approx\qty{9.9}{m/s}$。
\kotae （A）約\qty{9.0}{m/s}、（B）約\qty{6.5}{m/s}
\end{phreidai}
\begin{jibun}
質量\qty{20}{kg}のそりが、高さ\qty{3.0}{m}の所から速さ0で、角 $30^\circ$ の坂と、角 $15^\circ$ の坂をすべり下ります。動摩擦係数0.20として、下に着いたときの速さを、それぞれ求めましょう。
\bkjans{$30^\circ$：約\qty{6.2}{m/s}、$15^\circ$：約\qty{3.9}{m/s}}
\end{jibun}

\label{ph:力-4:ex13}
\begin{phreidai}[位置エネルギーが作れる力、作れない力]
平面の上の2つの力 $\pxev F_1=(y,\ x)$ と $\pxev F_2=(-y,\ x)$ について、点 $(0,\ 0)$ から点 $(1,\ 1)$ まで、（A）$(1,\ 0)$ を通る道と、（B）$(0,\ 1)$ を通る道の、2通りで仕事を計算しましょう。また、保存力かどうかを調べましょう（力の単位はN、長さはm）。
\phbamen
保存力なら、道によらず同じ仕事になるはずです。それを確かめます。
\phmushi
力は場所だけで決まり、時間で変わらないとします。
\phzu
\noindent\begin{minipage}[c]{0.36\linewidth}\centering
\begin{tikzpicture}[x=1mm, y=1mm]
  \draw[bkGray, line width=0.6pt, -{Stealth[length=3pt]}] (-2,0) -- (30,0) node[anchor=west, font=\sffamily\footnotesize] {$x$};
  \draw[bkGray, line width=0.6pt, -{Stealth[length=3pt]}] (0,-2) -- (0,30) node[anchor=south, font=\sffamily\footnotesize] {$y$};
  \draw[bk arrow, bkBlue, line width=1.1pt] (0,0) -- (24,0);
  \draw[bk arrow, bkBlue, line width=1.1pt] (24,0) -- (24,24);
  \draw[bk arrow, bkOrange, line width=1.1pt] (0,0) -- (0,24);
  \draw[bk arrow, bkOrange, line width=1.1pt] (0,24) -- (24,24);
  \node[font=\sffamily\footnotesize, text=bkBlue!80!black, anchor=north] at (14,-1) {A};
  \node[font=\sffamily\footnotesize, text=bkOrange!75!black, anchor=south] at (10,25) {B};
  \fill[bkInk] (24,24) circle (0.8);
  \node[font=\sffamily\footnotesize, anchor=west] at (25,25) {$(1,1)$};
\end{tikzpicture}
\end{minipage}\hfill
\begin{minipage}[c]{0.6\linewidth}\raggedright
道Aは、$x$ 軸にそって $(1,\ 0)$ まで進み、そこから真上に $(1,\ 1)$ まで。道Bは、$y$ 軸にそって $(0,\ 1)$ まで進み、そこから右へ $(1,\ 1)$ まで。\par\vspace{2pt}
横に進む所では $\pxev F\cdot\dd\pxev r=F_x\,\dd x$、たてに進む所では $F_y\,\dd y$ です。
\end{minipage}
\phriyu
線積分（④）を、まっすぐな2つの区間に分けて計算します。保存力かどうかは、$\dfrac{\partial F_y}{\partial x}-\dfrac{\partial F_x}{\partial y}$（⑥）で調べます。
\phtateru
$\pxev F_1$：A $=\displaystyle\int_0^1 0\,\dd x+\int_0^1 1\,\dd y$、B $=\displaystyle\int_0^1 0\,\dd y+\int_0^1 1\,\dd x$\\
$\pxev F_2$：A $=\displaystyle\int_0^1 0\,\dd x+\int_0^1 1\,\dd y$、B $=\displaystyle\int_0^1 0\,\dd y+\int_0^1(-1)\,\dd x$
\phtoku
（Aの横の区間は $y=0$、たての区間は $x=1$。Bのたての区間は $x=0$、横の区間は $y=1$。）\\
$\pxev F_1$：A $=0+1=\qty{1}{J}$、B $=0+1=\qty{1}{J}$。$\dfrac{\partial F_y}{\partial x}-\dfrac{\partial F_x}{\partial y}=1-1=0$。\\
$\pxev F_2$：A $=0+1=\qty{1}{J}$、B $=0+(-1)=\qty{-1}{J}$。$\dfrac{\partial F_y}{\partial x}-\dfrac{\partial F_x}{\partial y}=1-(-1)=2$。
\phtashikame
$\pxev F_1$ は $U=-xy$ で、$U(0,\ 0)-U(1,\ 1)=0-(-1)=1$ と合います（保存力）。\\
$\pxev F_2$ は道で仕事が変わるので保存力ではありません。Aで行ってBでもどる1周の仕事は $1-(-1)=2$ で、「rot の値2×囲んだ面積1」と同じです（8-11 のストークスの定理）。
\kotae $\pxev F_1$：どちらの道も\qty{1}{J}、保存力。$\pxev F_2$：Aは\qty{1}{J}、Bは\qty{-1}{J}、保存力ではない
\end{phreidai}
\begin{jibun}
$\pxev F_3=(2x,\ 2y)$ と $\pxev F_4=(y,\ -x)$ について、例題13と同じ2つの道で、$(0,\ 0)$ から $(1,\ 1)$ までの仕事を求め、保存力かどうかを答えましょう。
\bkjans{$\pxev F_3$：A・Bとも\qty{2}{J}、保存力（$\dfrac{\partial F_y}{\partial x}-\dfrac{\partial F_x}{\partial y}=0-0=0$、$U=-(x^2+y^2)$）。$\pxev F_4$：Aは\qty{-1}{J}、Bは\qty{1}{J}、保存力ではない}
\end{jibun}

\label{ph:力-4:ex14}
\begin{phreidai}[ばねにつけた物を、摩擦のある床ではなす]
水平な床の上で、ばね定数\qty{50}{N/m}のばねに質量\qty{0.50}{kg}の物体をつなぎ、ばねを自然の長さから\qty{0.10}{m}のばして、そっとはなしました。床との動摩擦係数を0.25とします。物体がはじめて止まる位置を求めましょう。
\phbamen
ばねに引きもどされて動き出し、自然の長さを通りすぎて、ばねを縮めてから止まります。摩擦で、行きすぎる距離は小さくなります。
\phmushi
ばねの質量と空気の抵抗を無視します。静止摩擦係数も0.25とします。
\phzu
自然の長さの位置を $x=0$ とし、はなした位置を $x_0=\qty{0.10}{m}$、はじめて止まる位置を $x=-a$（反対側に $a$）とします。動摩擦力の大きさ $\mu'mg$。
\phriyu
ばねの力は保存力、摩擦力は保存力ではないので、⑦の「力学的エネルギーの変化＝保存力以外の力の仕事」を使います。はなしたときも止まったときも速さ0なので、運動エネルギーは入りません。
\phtateru
動いた距離は $x_0+a$ です。
\[\frac{1}{2}ka^2-\frac{1}{2}kx_0^{\,2}=-\mu'mg\,(x_0+a)\]
\phtoku
左辺を $\dfrac{1}{2}k(a-x_0)(a+x_0)$ と因数分解し、両辺を $(x_0+a)$ でわると、
\[\frac{1}{2}k(a-x_0)=-\mu'mg\qquad a=x_0-\frac{2\mu'mg}{k}\]
\[a=0.10-\frac{2\times0.25\times0.50\times9.8}{50}=0.10-0.049=\qty{0.051}{m}\]
\phtashikame
大きさ：摩擦がなければ $a=x_0$（反対側の同じ所まで行く）。摩擦の分、行きすぎが $\dfrac{2\mu'mg}{k}=\qty{0.049}{m}$ だけ短くなりました。\\
止まったままか：$x=\qty{-0.051}{m}$ でのばねの力 $50\times0.051\approx\qty{2.6}{N}$ は、最大の静止摩擦力 $0.25\times0.50\times9.8\approx\qty{1.2}{N}$ より大きいので、物体はまた動き出します。
\kotae 自然の長さから反対側へ約\qty{5.1}{cm}の所
\end{phreidai}
\begin{jibun}
ばね定数\qty{200}{N/m}のばねに、質量\qty{1.0}{kg}の物体をつなぎ、\qty{0.15}{m}のばしてはなしました。動摩擦係数0.10として、はじめて止まる位置を求めましょう。
\bkjans{$a=0.15-0.0098$ より、反対側へ約\qty{0.14}{m}}
\end{jibun}

\label{ph:力-4:ex15}
\begin{phreidai}[半球の上をすべり下りる]
半径\qty{1.2}{m}のなめらかな半球（ドーム）のいちばん上に、小さな物体を置き、ほんの少し押したら、すべり下り始めました。物体が半球の面から離れる高さ（半球の中心からの高さ）と、そのときの速さを求めましょう。
\phbamen
はじめは面にそって円運動をしますが、速くなると、面が物体を支えられなくなり、面から離れて飛び出します。
\phmushi
摩擦・空気の抵抗を無視し、物体は点とします。いちばん上での速さは0とします。
\phzu
\noindent\begin{minipage}[c]{0.4\linewidth}\centering
\begin{tikzpicture}[x=0.82mm, y=0.82mm]
  \fill[bkGray!25] (-24,0) arc (180:0:24) -- cycle;
  \draw[bkInk, line width=1pt] (-24,0) arc (180:0:24);
  \draw[bkInk, line width=0.6pt] (-28,0) -- (28,0);
  \draw[bkInk, line width=0.5pt] (0,0) -- (0,24);
  \draw[bkInk, line width=0.5pt] (0,0) -- (48.19:24);
  \draw[bkInk, line width=0.4pt] (0,7) arc (90:48.19:7);
  \node[font=\sffamily\footnotesize] at (5,10) {$\theta$};
  \fill[bkRed] (48.19:24) circle (1.2);
  \draw[bk arrow, bkPurple, line width=1pt] (48.19:24) -- ++(0,-9);
  \node[font=\sffamily\footnotesize, text=bkPurple!80!black, anchor=west] at ($(48.19:24)+(0.8,-7)$) {$mg$};
  \draw[bkGray, dashed, line width=0.5pt] (48.19:24) -- (30,17.9);
  \node[font=\sffamily\footnotesize, anchor=north] at (-6,-1) {中心};
\end{tikzpicture}
\end{minipage}\hfill
\begin{minipage}[c]{0.57\linewidth}\raggedright
いちばん上から角 $\theta$ だけ回った所で、速さ $v$、面が押す力 $N$（中心から外向き）とします。そこの高さは中心から $r\cos\theta$ です（$r=\qty{1.2}{m}$）。重力の、中心向きの成分は $mg\cos\theta$ です。
\end{minipage}
\phriyu
例題2と同じく、「離れる条件 $N=0$」は円運動の運動方程式（中心向き）で、速さは力学的エネルギー保存で求めます。
\phtateru
\[m\frac{v^2}{r}=mg\cos\theta-N,\qquad \frac{1}{2}mv^2=mgr(1-\cos\theta)\]
\phtoku
2つめの式から $v^2=2gr(1-\cos\theta)$。これを1つめの式に入れて、
\[N=mg\cos\theta-2mg(1-\cos\theta)=mg(3\cos\theta-2)\]
$N=0$ になるのは、$\cos\theta=\dfrac{2}{3}$ のときです。
\[\text{高さ}=r\cos\theta=\frac{2}{3}\times1.2=\qty{0.80}{m}\]
\[v=\sqrt{2gr\times\frac{1}{3}}=\sqrt{\frac{2\times9.8\times1.2}{3}}=\qty{2.8}{m/s}\]
\phtashikame
角：$\cos\theta=\dfrac{2}{3}$ から $\theta\approx48^\circ$（電卓の $\cos^{-1}$）。\\
半径によらない：離れる高さは、いつも半径の $\dfrac{2}{3}$ です（$g$ にも $m$ にもよらない）。
\kotae 中心から\qty{0.80}{m}の高さ、速さ\qty{2.8}{m/s}
\end{phreidai}
\begin{jibun}
半径\qty{0.60}{m}のなめらかな半球のいちばん上から、物体がすべり下りました。面から離れる高さ（中心から）と、そのときの速さを求めましょう。
\bkjans{高さ\qty{0.40}{m}。速さは $v^2=2\times9.8\times0.60\div3$ より約\qty{2.0}{m/s}}
\end{jibun}

%% ======================== まちがえやすいところ ========================
\Needspace{18\baselineskip}\bksection{tsumazuki}{まちがえやすいところをチェック}
\begin{tsumazuki}{速さが2倍なら、エネルギーも2倍？}
\bkng{時速\qty{40}{km}から\qty{80}{km}にすると、運動エネルギーも止まるまでの距離も2倍。}
\bkok{運動エネルギーは $v^2$ に比例するので4倍。止まるまでの距離も4倍。}
$K=\dfrac{1}{2}mv^2$ の2乗を忘れると、ブレーキの距離を大きく見誤ります（⑤）。
\end{tsumazuki}

\begin{tsumazuki}{垂直抗力も仕事をする？}
\bkng{坂を下りるとき、垂直抗力も、重力と同じように仕事をする。}
\bkok{垂直抗力は、動く向きと直角なので、仕事は0。}
仕事は「力×力の向きに動いた距離」。直角な力は、どんなに大きくても仕事をしません（④）。
\end{tsumazuki}

\begin{tsumazuki}{効率をかけるか、わるか}
\bkng{効率0.60のポンプに入れる動力は、有効な動力×0.60。}
\bkok{入れる動力＝有効な動力÷0.60。入れるほうがいつも大きい。}
効率は「出たもの÷入れたもの」。答えが「入れた動力より大きな有効な動力」になったら、かけ算とわり算が逆です（例題9）。
\end{tsumazuki}

%% ======================== 工学ではこう使う ========================
\Needspace{18\baselineskip}\bksection{phgear}{工学ではこう使う}
\begin{phkougaku}[動力で機械を選ぶ]
モーターやエンジンを選ぶときは、まず「1秒あたりにどれだけの仕事をさせたいか」（動力）を計算します。荷物を上げるなら $mgv$、ポンプなら $\rho QgH$、車なら（空気の抵抗などを無視した坂の分だけでも）$mg\sin\theta\times v$ です。これを効率でわったものが、機械に入れる動力です。

電車やハイブリッド車・電気自動車のブレーキの一部は、モーターを発電機として使い、運動エネルギーを電気にもどしてためます（回生ブレーキ）。摩擦のブレーキでは熱になって失われるエネルギーを、もう一度使うしくみです。
\par\smallskip\noindent\begin{center}
\begin{tikzpicture}[x=1mm, y=1mm]
  \node[bk node, minimum width=30mm, minimum height=10mm, font=\sffamily\small] (k) at (0,0) {運動エネルギー\\$\tfrac{1}{2}mv^2$};
  \node[bk node, minimum width=30mm, minimum height=10mm, font=\sffamily\small] (e) at (66,0) {電気のエネルギー\\（電池にためる）};
  \node[bk node, minimum width=30mm, minimum height=8mm, font=\sffamily\small, draw=bkRed] (h) at (0,-24) {熱（摩擦のブレーキ）};
  \draw[bk link] ([yshift=2mm]k.east) -- node[above=1mm, font=\sffamily\footnotesize, text=bkInk] {回生ブレーキ} ([yshift=2mm]e.west);
  \draw[bk link] ([yshift=-2mm]e.west) -- node[below=1mm, font=\sffamily\footnotesize, text=bkInk] {モーターで走る} ([yshift=-2mm]k.east);
  \draw[bk link, bkRed] (k.south) -- (h.north);
\end{tikzpicture}
\end{center}

\begin{phketa}[仕事率とエネルギー]
\phketaline{人が8時間の仕事で続けられる仕事率}{約\qty{75}{W}}
\phketaline{軽自動車のエンジンの最高出力（上限）}{\qty{47}{kW}（64馬力）}
\phketaline{質量\qty{1000}{kg}の車が時速\qty{100}{km}で走る運動エネルギー}{約\qty{0.39}{MJ}}
\phketaline{水力発電の出力（kW）}{約 $9.8\times Q\times H\times\eta$}
\phboxhead{答えが桁で合っているか}
人の力でできる仕事は数十〜数百 W、自動車は数十〜百 kW、大きな発電所は数百 MW〜GW（\qty{1}{GW}＝\qty{1e9}{W}）です。人が数 kW を出し続ける答えが出たら、どこかで単位（g と kg、L と m$^3$、分と秒）をまちがえています。
\end{phketa}
\end{phkougaku}
\phstation{動力伝達,流体機械,機械力学}{仕事率、動力、効率、揚程、ポテンシャル、保存力、回生}

\begin{column}{ワットと馬力}
仕事率の単位 W（ワット）は、蒸気機関を大きく改良したイギリスのジェームズ・ワット（1736〜1819年）の名前からとりました。ワットは、自分の蒸気機関の力を、それまで使われていた馬とくらべて売りこむために、「1頭の馬が続けてできる仕事率」として\textbf{馬力}を決めました。この英馬力は約\qty{746}{W}です。日本で車のエンジンの説明に使われてきた馬力（仏馬力、PS）は、約\qty{735.5}{W}で、少しだけちがいます。今の計量の決まりでは、仕事率は W で表します。
\end{column}

\begin{column}{ガリレオの「くぎを打ったふりこ」}
\bkyattemiyouref のふりこの実験は、ガリレオ・ガリレイ（1564〜1642年）が1638年に出版した本『新科学対話』の中で説明しているものです。かべにくぎを打ち、ふりこの糸がくぎに引っかかるようにしても、おもりは放した高さとほぼ同じ高さまで上がります。ガリレオは、これを使って「同じ高さから下りてきた物は、どんな道を通っても、同じ高さまで上がれるだけの勢いをもつ」と考えました。まだ「エネルギー」ということばがなかった時代に、この章の力学的エネルギー保存の考えにたどり着いていたのです。エネルギーということばが今の意味で使われるようになったのは、それから200年ほどあとの19世紀のことです。
\end{column}

%% ======================== 章末問題 ========================
\bkendproblems
\bkproblemgroup{基本}{3}
\bkproblem 高さ\qty{45}{m}の山の上から速さ0で下り始めたコースターが、高さ\qty{5.0}{m}の所を通るときの速さを求めましょう（摩擦なし）。
\bkpans{$\sqrt{2\times9.8\times40}=\qty{28}{m/s}$}
\bkproblem 質量\qty{25}{kg}の荷物を、角 $20^\circ$ のなめらかな斜面にそって、ゆっくり\qty{5.0}{m}押し上げます。押す力と仕事を求めましょう（$\sin20^\circ\approx0.342$）。
\bkpans{力 約\qty{84}{N}、仕事 約\qty{4.2e2}{J}}
\bkproblem ばね定数\qty{300}{N/m}のばねを\qty{6.0}{cm}縮めて、質量\qty{30}{g}の玉を水平に飛ばします。飛び出す速さを求めましょう。
\bkpans{$0.060\times\sqrt{300\div0.030}=\qty{6.0}{m/s}$}
\bkproblem 速さ\qty{15}{m/s}の車が急ブレーキですべって止まりました。動摩擦係数0.50として、すべった距離を求めましょう。
\bkpans{$15^2\div(2\times0.50\times9.8)\approx\qty{23}{m}$}
\bkproblem クレーンが、質量\qty{500}{kg}の荷物を一定の速さ\qty{0.80}{m/s}で引き上げます。ワイヤが引く仕事率を求めましょう。
\bkpans{$500\times9.8\times0.80=\qty{3.92e3}{W}$（約\qty{3.9}{kW}）}
\bkproblem 毎分\qty{120}{L}の水を\qty{8.0}{m}の高さまでくみ上げます。有効な動力と、効率0.70のときに入れる動力を求めましょう。
\bkpans{有効な動力 約\qty{157}{W}、入れる動力 約\qty{224}{W}}
\bkproblemgroup{チャレンジ}{4}
\bkproblem 半径\qty{6.0}{m}のループを摩擦なしで1周するのに必要な、いちばん低い出発の高さを求めましょう。そのとき、いちばん下で座席が人を押す力は体重の何倍ですか。
\bkpans{\qty{15}{m}。6倍}
\bkproblem 火星の $GM=\qty{4.283e13}{m^3/s^2}$、半径 $\qty{3.390e6}{m}$ として、火星の表面から飛び出す速さを求めましょう。
\bkpans{$v^2=2\times\num{4.283e13}\div\num{3.390e6}$ より、約\qty{5.0}{km/s}}
\bkproblem 質量\qty{20}{kg}のそりが、高さ\qty{3.0}{m}の所から速さ0で、角 $45^\circ$ の坂と角 $20^\circ$ の坂をすべり下ります。動摩擦係数0.20として、下に着いたときの速さを、それぞれ求めましょう（$\tan20^\circ\approx0.364$）。
\bkpans{$45^\circ$：約\qty{6.9}{m/s}、$20^\circ$：約\qty{5.1}{m/s}}
\bkproblem 半径\qty{0.90}{m}のなめらかな半球のいちばん上から、物体がすべり下りました。面から離れる高さ（中心から）と、そのときの速さを求めましょう。
\bkpans{\qty{0.60}{m}、約\qty{2.4}{m/s}}
\bkproblem 人と自転車で質量\qty{60}{kg}が、こう配6\,\%の坂を一定の速さ\qty{5.0}{m/s}で上ります。坂を上るためだけの仕事率を求めましょう（$\sin\theta\approx\tan\theta$ としてよい）。
\bkpans{$60\times9.8\times0.06\times5.0$ より、約\qty{1.8e2}{W}}
\bkproblem 毎秒\qty{5.0}{m^3}の水が、高さの差\qty{40}{m}を下りる発電所の出力を、効率0.80として求めましょう。
\bkpans{$0.80\times1000\times5.0\times9.8\times40$ より、約\qty{1.6}{MW}}
\bkproblem 平面の力 $\pxev F=(2xy,\ x^2)$ が保存力かどうかを調べましょう。また、例題13の2つの道で、点 $(0,\ 0)$ から点 $(1,\ 1)$ までの仕事を求めましょう（力はN、長さはm）。
\bkpans{$\dfrac{\partial F_y}{\partial x}-\dfrac{\partial F_x}{\partial y}=2x-2x=0$ で保存力（$U=-x^2y$）。仕事はどちらの道も\qty{1}{J}}
\phnyushi
\bkproblem 鉛直に立てた半径 $r$ のなめらかな円筒の内側で、いちばん下にある小球に、内面にそって水平に速さ $v_0$（$v_0^2=3gr$）を与えました。重力加速度を $g$、小球の質量を $m$ とします。（1）いちばん下で面が小球を押す力、（2）小球が面から離れる位置の、いちばん下からの高さ、（3）そのときの小球の速さを、文字で答えましょう。
\bkpans{（1）\nobreak $4mg$。（2）\nobreak $\dfrac{4}{3}r$（中心より $\dfrac{r}{3}$ 上）。（3）\nobreak $\sqrt{\dfrac{gr}{3}}$}
\bkproblem 水平な床の上で、ばね定数\qty{100}{N/m}のばねに質量\qty{1.0}{kg}の物体をつなぎ、自然の長さから\qty{0.20}{m}のばしてはなしました。静止摩擦係数・動摩擦係数はどちらも0.20とします。（1）物体がはじめて止まる位置、（2）物体が最後に止まる位置、（3）止まるまでに動いた道のりの合計を求めましょう。
\bkpans{（1）\nobreak 反対側へ\qty{0.16}{m}。（2）\nobreak 止まるたびに\qty{0.0392}{m}ずつ近づき、5回めに反対側の\qty{0.004}{m}（\qty{4}{mm}）で止まったまま。（3）\nobreak ばねのエネルギーの減った分 約\qty{2.0}{J}を、摩擦力\qty{1.96}{N}でわって、約\qty{1.02}{m}}

\begin{furikaeri}
\begin{checklist}
\item 高さと速さの関係を、仕事と運動エネルギーで説明できる。\bkfurref{ph:力-4:sec3}
\item エネルギーの保存と摩擦の損失から、速さ・高さ・仕事率を求められる。\bkfurref{ph:力-4:sec7}
\item $\pxev F=-\mathrm{grad}\,U$ から、位置エネルギーの式を導ける。\bkfurref{ph:力-4:sec6}
\end{checklist}
\end{furikaeri}
\bkfukidashi{mina}{高さは「速さの貯金」なんだね。摩擦があると貯金が少しずつ減っていくから、最初の山より高くは上れないんだ。}
